% !TeX program = pdflatex
% ------------------------------------------------------------------
% arXiv submission draft.  Suggested classes:
%   primary: math.AT ; secondary: cs.RO, cs.DC
% ------------------------------------------------------------------
\documentclass[12pt]{article}

\setlength{\oddsidemargin}{.5cm}
\setlength{\evensidemargin}{.5cm}
\setlength{\textwidth}{6in}
\setlength{\textheight}{8in}
\setlength{\emergencystretch}{1.5em}

\usepackage{amsmath,amssymb,amsthm}
\usepackage{array,longtable}
\usepackage{float}
\usepackage{algorithm}
\usepackage{algpseudocode}
\usepackage[font=small,labelfont=bf]{caption}
\usepackage{booktabs}
\usepackage{hyperref}
\hypersetup{colorlinks=true,linkcolor=blue,citecolor=blue,urlcolor=blue}

\newcommand{\doi}[1]{doi:\,\href{http://dx.doi.org/#1}{#1}}
\newcommand{\TC}{\mathrm{TC}}
\newcommand{\cat}{\mathrm{cat}}
\newcommand{\SC}{\mathrm{SC}}
\newcommand{\SCst}{\SC_{\mathrm{strict}}}
\newcommand{\SD}{\mathrm{SD}}

\newtheorem{theo}{Theorem}[section]
\newtheorem{lemma}[theo]{Lemma}
\newtheorem{corollary}[theo]{Corollary}
\newtheorem{proposition}[theo]{Proposition}
\newtheorem{definition}[theo]{Definition}
\newtheorem{example}[theo]{Example}
\newtheorem{remark}[theo]{Remark}
\newtheorem{conjecture}[theo]{Conjecture}

\title{Explicit covers for simplicial complexity:\\
 from spheres to surfaces.}

\author{M. Perez\footnote{Escuela Superior de F\'isica
y Matem\'aticas del Instituto Polit\'ecnico Nacional, Edificio 9, Unidad
Profesional Adolfo L\'opez Mateos, 07300 Mexico City, Mexico.}, and
J. Flores\footnotemark[1]}

\date{}

\begin{document}
\maketitle

\begin{abstract}
Computing Gonz\'alez's strict simplicial complexity $\SCst(K)$ at
subdivision level zero requires an explicit cover of $K\times K$ by
contiguity domains.  In the authors' own $96$-core implementation of the
randomized greedy pipeline of \cite{OLGB21}, $120$-hour jobs never
completed the $2$-sphere model.  We diagnose the failure and describe
the three-stage pipeline that replaced it.  (i)~The \emph{star cover}.
For connected $K$ every star rectangle $D(v,w)$, generated by the product
facets over $\operatorname{star}(v)\times\operatorname{star}(w)$, is a
contiguity domain; the exact minimum number $\rho(K)$ of rectangles
covering $K\times K$ is a small set cover, so in seconds and without
search one obtains a complete verified cover, the bound
$\SCst(K)\le\rho(K)-1$, and chains that seed the next stage.  The bound
is sharp on every even sphere $\partial\Delta^{n+1}$, where
$\SCst=\TC(S^n)=2$.  (ii)~A \emph{chain-first annealer} that evolves
whole systems of $k$ contiguity chains on a single GPU, warm-started from
the star cover and descended in $k$ below $\rho$.  It finds three-domain
systems for $S^2$ through $S^6$ (the cube boundary, with $\rho=8$,
included), and on the icosahedron boundary descends from the
$24$-domain star cover to a verified $5$-domain cover, so that
$2\le\SCst\le4$, where $k=4$ hard-stalls at $2376/2400$ facets.
(iii)~\emph{CP-SAT} decides small instances exactly.  It shows that no
two-domain cover of $\partial\Delta^4\times\partial\Delta^4$ with chains
of at most $48$ links exists, where the annealer's $k=2$ searches stall at
$460/500$.  This is a conditional solver result: under a compression
conjecture on chain lengths, which we state and do not prove, it makes
$\partial\Delta^4$ a second example, after the M\"obius band, of a
level-zero gap $\SCst>\TC$.  Every reported cover is verified link by
link, and all but three families are re-checked by an independent
auditor.
\end{abstract}

\medskip\noindent
\emph{Keywords and phrases:} topological complexity; simplicial
complexity; contiguity chains; piecewise linear motion planning;
simulated annealing; GPU computing.

\medskip\noindent
\emph{2020 Mathematics Subject Classification:} Primary: 55M30, 68W20.
Secondary: 68T40, 55U10, 90C27.

\section{Introduction}\label{sec:intro}

Farber's topological complexity $\TC(X)$ \cite{F03} measures the minimum
number of continuous motion-planning rules on a space $X$.
We use the reduced convention throughout: $\TC(X)+1$ is the minimum
number of local motion-planning rules, and a cover by $k$ such
domains certifies $\TC(X)\le k-1$.
Gonz\'alez's simplicial complexity \cite{G18} discretizes it: on the
ordered product $K\times K$ of a triangulation $K$ of $X$, continuity
becomes contiguity of simplicial maps (Section~\ref{sec:back}).  The
resulting invariant $\SCst(K)$ is finite and checkable, but computing
it requires an explicit cover of $K\times K$ by contiguity domains.

The pipeline of \cite{OLGB21} attacks this problem by growing one
domain at a time via budgeted random walks
(Section~\ref{sec:back}).  Their published experiments reach the circle
($18$ product facets, ${\approx}1$\,min) and the wedge of two circles
($50$ facets, $25$\,min); the $2$-sphere was deferred to cluster
hardware.  A C++/MPI predecessor of the present work ran the same
algorithms on $96$ cores in $120$-hour jobs (runs by the authors with
their own implementation, not measurements reported in
\cite{OLGB21}) and never completed one
covering of $\partial\Delta^3\times\partial\Delta^3$
(Section~\ref{sec:fail}).  A faithful single-GPU port did produce a
cover of that model, in $4.4$\,minutes.
The initial GPU baseline produced a $20$-domain cover of the $3$-sphere
model.  A later run under a batched-retry configuration produced an
independently checked $13$-domain cover, certifying
$\SCst(\partial\Delta^4)\leq12$; it does not improve the three-domain
chain-first certificate.

This paper reports what we did next, in the order we did it, and the
pipeline we now recommend.  Its spine is a three-stage attack on the
minimum cover of $K\times K$.

\begin{enumerate}
\item[(i)] The \textbf{star cover}: build the $|V|^2$ product rectangles
$D(v,w)$, each a contiguity domain by Theorem~\ref{thm:prod}, and solve
the minimum subcover $\rho(K)$ exactly.  It takes seconds on a CPU, needs
no search, and returns a complete verified cover, the unconditional bound
$\SCst(K)\le\rho(K)-1$, and chains that seed stage (ii).
\item[(ii)] The \textbf{chain-first annealer}: anneal whole systems of
$k$ contiguity chains on a GPU, warm-started from the star chains and
descended in $k$ below $\rho$, one unit at a time.
\item[(iii)] \textbf{CP-SAT}: when the annealer hard-stalls at a $k$
above the topological floor, an exact solver can sometimes decide the
instance.  A stall is evidence, not proof, that no $k$-cover exists;
CP-SAT can settle it at small scale (the M\"obius band, the $3$-sphere
model), and it does not scale beyond.
\end{enumerate}

\noindent We tell the story in the order we discovered the pieces,
because the order is part of the argument.  First came the annealer alone,
which we ran from a large $k$ where a first cover is easy and descended
(Section~\ref{sec:descent}); the cost of that first cover, which was
hours on the larger surfaces, motivated everything after.  Only later
did we realize that a theorem suggested by the annealer's own plateaus is
a practical first move: the star cover.  There are five contributions,
in the order the reader meets them.

\begin{itemize}
\item \textbf{Diagnosis} (Section~\ref{sec:fail}).  We identify two
mechanisms consistent with the recorded stalls of the greedy pipeline:
an exhausted extension call returns only failure, without grading how
close the attempt came, while contiguity acts as a hard constraint; and
sequential, irrevocable growth leaves near-product fronts at which no
tested extension succeeded within the recorded budgets.
\item \textbf{A chain-first annealer} (Section~\ref{sec:annealer}).
Instead of growing domains facet by facet, we anneal whole systems of
$k$ complete contiguity chains in parallel on one GPU, under an energy
that grades every uncovered facet by the defect count of its best
chain and penalises defects in all chains at once.  The search finds explicit covers:
a two-domain system for $S^1$ and three-domain systems for $S^2$ through $S^6$---$22$ seconds
for the $2$-sphere, ${\approx}2.5$ hours for the $6$-sphere---certifying
$\SCst(\partial\Delta^{n+1})\le2$ for all $n\le6$, sharply ($=\TC$) on
$S^1$ (where the certified bound is $\SCst\le1$), $S^2$, $S^4$, $S^6$ and the cube boundary (Table~\ref{tab:annealer}).
\item \textbf{The product-domains theorem and the star cover}
(Section~\ref{sec:products}).  The two-chain plateau runs we recorded
stopped at coverage values equal to those attained by two star
rectangles $D(v,w)$.  We prove that, for connected $K$, every such
rectangle is a contiguity domain (Theorem~\ref{thm:prod}), giving
$\SCst(K)\le\rho(K)-1$ for the exactly computable rectangle-cover
constant $\rho(K)$ (Corollary~\ref{cor:rho}); for the simplex boundary
$\rho(\partial\Delta^{n+1})=3$, hence
\begin{gather*}
\SCst(\partial\Delta^{n+1})\le 2\ \ \text{for all }n\ge1,\\
\SCst(\partial\Delta^{n+1})=\TC(S^n)=2\ \ \text{for even }n\ge2,
\end{gather*}
sharp at subdivision level $0$ (Corollary~\ref{cor:spheres}).  A
counting lemma (Lemma~\ref{lem:rows}) brackets $\rho$ and computes it
outright on several complexes, and we tabulate $\rho$ with its bounds for
the complexes of our campaigns (Table~\ref{tab:rho}).  Presented as a method, the star cover
is the canonical first move and a seed (Section~\ref{sec:starmethod}).
\item \textbf{Descents, seeds and a case study}
(Sections~\ref{sec:descent} and~\ref{sec:icosa}).  We document the
$k$-descent ladder, the cost of the cold first cover, and the star-seeded
descents on the surfaces.  On the icosahedron boundary the star cover
($24$ domains) seeds a descent to a verified $5$-domain cover,
$\SCst\le4$, nineteen domains below its seed and within two of the
topological floor; at $k=4$ three runs hard-stall at $2376/2400$.
\item \textbf{Exact decisions} (Section~\ref{sec:sat}).  A validated
CP-SAT encoding shows that no $2$-domain cover of
$\partial\Delta^4\times\partial\Delta^4$ with chains of at most $48$
links exists.  This is a solver result for a bounded class of covers; it
settles Conjecture~\ref{conj} (C1) within that class, and upgrades it
to $\SCst(\partial\Delta^4)=2$ only given a compression conjecture
(Conjecture~\ref{conj:compress}) that we measure but do not prove.
\item \textbf{Certificates, not claims} (Section~\ref{sec:repro}).
Every reported system is replayed deterministically and re-audited by
an independent checker sharing no code with the search; the eight
archived certificates of the first phase are checked by a pinned
auditor against a hash manifest, the star-cover and ladder certificates
of the surface campaigns were re-audited for this paper except for three
complexes the auditor cannot yet rebuild, and source and certificates
are public.
\end{itemize}

The stages are complementary, and
Sections~\ref{sec:results} and~\ref{sec:concl} keep them in tension.  On
the cube boundary, where covers by the star rectangles $D(v,w)$ of
Theorem~\ref{thm:prod} need $\rho=8$ domains and certify only
$\SCst\le7$, the annealer finds three domains, which is optimal since
$\TC(S^2)=2$.  On the dodecahedron boundary a $20$-hour ladder campaign
at $k=3$ finds no complete cover within its budget, and the star cover
remains the only certified bound ($\SCst\le35$).  The exact level-$0$
value on $\partial\Delta^4$ remains open pending a compression claim.
The results, old and new, are collected in Table~\ref{tab:records}
and Table~\ref{tab:surfaces}.

\section{Background and notation}\label{sec:back}

We keep the framework of \cite{OLGB21} and recall only what is used
later.  Throughout, $K$ is a finite, nonempty, connected simplicial complex on an
ordered vertex set $V$, and $K\times K$ is the ordered (``staircase'')
product triangulation of $\|K\|\times\|K\|$; its maximal facets lie
over pairs $(s,t)$ of facets of $K$.

Two simplicial maps $\varphi,\varphi':J\to K$ are \emph{$1$-contiguous}
if $\varphi(\sigma)\cup\varphi'(\sigma)$ is a simplex of $K$ for every
facet $\sigma$ of $J$; a \emph{contiguity chain} is a sequence
$\varphi=\varphi_0\sim\varphi_1\sim\cdots\sim\varphi_c=\varphi'$ of
consecutively $1$-contiguous maps.  A subcomplex $J\subseteq L$ is a
\emph{contiguity domain} for $\psi,\psi':L\to K$ if the restrictions
$\psi|_J,\psi'|_J$ are joined by a contiguity chain, and the strict
contiguity distance $\SD(\psi,\psi')$ is one less than the minimum
number of contiguity domains covering $L$.  For the two projections
$\pi_1,\pi_2:K\times K\to K$,
\[
\SCst(K)=\SD(\pi_1,\pi_2)\ge \TC(\|K\|),
\]
with equality in the limit of barycentric subdivisions
\cite{OLGB21,G18}.  The rules counted by $\TC$ are continuous local
sections of the path-evaluation map over an open cover of
$\|K\|\times\|K\|$, following \cite{F03}, whereas contiguity domains
are subcomplexes; the comparison between the two invariants is the
theorem of \cite{G18}, not a literal identification of domains with
open sets.  A cover of $K\times K$ by $d$ contiguity domains is
a \emph{system of $d$ piecewise-linear motion planners} and certifies
$\SCst(K)\le d-1$.  All statements in this paper are at subdivision
level $0$, and we read the facets of $K\times K$ as a covering problem
over the \emph{facet pairs} $(s,t)$.  When $K$ is pure of dimension
$d$, every pair carries the same number $\binom{2d}{d}$ of staircase
facets, and a facet count is a pair count up to that factor; when $K$
is not pure, the number of staircase facets over $(s,t)$ depends on
$\dim s$ and $\dim t$.  In either case a coverage figure counted in
facets becomes a count of covered \emph{pairs} only after checking
that each pair is covered in its entirety, since a domain may contain
some of the staircase facets over a pair without containing them all.

The pipeline of \cite{OLGB21} grows one domain at a time: starting from
a seed facet it repeatedly calls \textsc{LocalSearch}---a random walk
with budget $M$ that tries to extend a contiguity chain over one more
facet---then merges and prunes the resulting cover.  Its cost profile,
and the way it fails, is the subject of the next section.

\section{How the greedy dies}\label{sec:fail}

We ran the pipeline of \cite{OLGB21} on the $2$-sphere model
($\partial\Delta^3\times\partial\Delta^3$, $96$ product facets) using
our faithful GPU port under the original settings: walk budget
$M=40000$, unbounded retry passes, $16$ replicas across two
runs.  These are the authors' own runs (this work), not measurements
reported in \cite{OLGB21}.  We then instrumented every frozen front the
pipeline produced.\footnote{Repository:
\url{https://github.com/MiguelPerezI/Abstract_Motion_GPU}.}
Two mechanisms are consistent with the stall.
The figures in this section are measurements reported by the authors
from their instrumented runs; the scripts, snapshots and logs needed to
reproduce them are not yet available in the public repositories, so
they have not been independently reproduced.

\textbf{Binary extension feedback and hard feasibility.}
\textsc{LocalSearch} uses the graph-distance score
$d(f,\pi_2)=\sum_{v\in V(J)}d_K(f(v),\pi_2(v))$
to guide admissible steps. A proposal that violates contiguity is rejected
regardless of its distance improvement, and an exhausted extension call
returns failure without grading how close the attempt came to success.
This differs from the chain-first energy, which explicitly scores
contiguity violations in candidate vertex assignments.
(See \cite[\S4.1, Algorithm~1]{OLGB21}.)  The distance $d$ to the
target map and the number of links that violate contiguity on a facet
(the quantity $b_j(\sigma)$ of Section~\ref{sec:annealer}) are different
quantities: a vertex change that decreases $d$ need not be admissible.

\textbf{Self-cornering.}  Growth is sequential and irrevocable---one
facet at a time, never undone---so each accepted facet constrains the
next, and in the recorded replicas every domain ended as a \emph{near-rectangle}: most of a
maximal product rectangle $D(v,w)$ of Theorem~\ref{thm:prod}, plus a
few facets outside it.  Every replica froze this way, at $30$ to $54$
covered facets.  To probe these fronts, we tested all
$146$ pending candidates at three frozen fronts with the original walk
(up to $3\times10^{8}$ steps each) and with a soundness-checked
single-chain annealer miniature at the parameters of
Table~\ref{tab:params}.
No tested extension succeeded within the recorded budgets, using either
the original walk or the single-chain annealer. These experiments show
persistent stagnation under the tested search policies; they do not
establish that an extension is impossible or that the front is a local
optimum under all admissible moves.  The wedge-of-two-circles
reproduction of the first phase exhibited a similar signature at
smaller scale.

These two mechanisms interact.  A failed extension call returns no
graded information, so the retry law cannot distinguish a candidate
that cannot be attached from one that was merely unlucky within its
budget.  In the instrumented run, $65$ consecutive retry passes over
the full frozen front produced zero successes and took $82\%$ of the
wall time without further growth.

\textbf{Step cost.}  Each walk step rescans the domain and re-runs
shortest paths from scratch---measured on the original C++ sources,
about $77\,\mu$s per step---so one failed $M{=}40000$ walk costs
${\approx}3$\,seconds.  Multiply by $96$ cores and $120$ hours and the
cluster's output is a large number of stalls, not a large covering.

\emph{A correction.}  An earlier reading of ours attributed the freeze
to chain drift---entries of a chain receiving no signal and randomizing
where a neighbouring chain needs them.  The greedy rebuilds every chain
from $\pi_1$ and accepts only fully contiguous links, so there is
nothing to drift; drift is a disease of the \emph{annealer's} energy,
suppressed by the $\gamma$ term of Section~\ref{sec:annealer}.

The repair of the step cost is engineering (batched GPU walks; on one
Tesla T4 the kernel issues $10^{10}$ steps in $4.4$\,minutes).  The
response to the two mechanisms is the energy of the next section:
grading defects replaces binary extension feedback, and annealing whole
systems avoids irrevocable sequential growth.

\section{The chain-first annealer}\label{sec:annealer}

Instead of growing domains facet by facet, we anneal \emph{whole
candidate systems}.  A state is a tuple of $k$ chains.  A chain is a
sequence $C=(\varphi_0,\varphi_1,\dots,\varphi_{c_{\max}})$ of maps of
vertex sets with fixed endpoints $\varphi_0=\pi_1$ and
$\varphi_{c_{\max}}=\pi_2$; the interior entries are arbitrary vertex
assignments and are \emph{not} required to be simplicial maps
$K\times K\to K$.  In the implementation $c_{\max}$ counts links, so a
stored chain holds $c_{\max}+1$ maps.

A facet $\sigma$ of $K\times K$ is \emph{covered} by $C$ when every
link of $C$ satisfies the $1$-contiguity condition on $\sigma$, and the
\emph{domain} $\mathrm{Dom}(C)$ is the subcomplex generated by the
facets that $C$ covers.  Restricted to $\mathrm{Dom}(C)$ the entries
are simplicial and consecutive entries are $1$-contiguous.  Indeed, for
$c_{\max}\ge1$, each image $\varphi_i(\sigma)$ of a covered facet
$\sigma$ is contained in the union of the images under an adjacent
link, which is a simplex of $K$; closure under faces makes the
restricted maps simplicial and extends contiguity to every face of the
domain, so no global simpliciality of the candidates is required.
Hence $C|_{\mathrm{Dom}(C)}$ is a contiguity chain from
$\pi_1|_{\mathrm{Dom}(C)}$ to $\pi_2|_{\mathrm{Dom}(C)}$ in the sense
of Section~\ref{sec:back}.  It is this restriction, not the ambient
vertex assignment, that a certificate reports.

For a facet $\sigma$ and the $j$-th chain of a state, let $b_j(\sigma)$
be the number of links of that chain which violate the contiguity
condition on $\sigma$, and let $m(\sigma)=\min_j b_j(\sigma)$ be the
defect count of the best chain for that facet.  With $\mathrm{pen}(0)=0$
and $\mathrm{pen}(b)=1+\beta b$ for $b>0$, where $\beta\ge0$,
$\gamma\ge0$ and $w_\sigma>0$, the energy of a state is
\[
E \;=\; \sum_{\sigma} w_\sigma\,\mathrm{pen}\bigl(m(\sigma)\bigr)
\;+\; \gamma\sum_{j}\sum_{\sigma} b_j(\sigma).
\]
The facet weights are $w_\sigma=1$ in the unguided runs; a guided
restart may raise the weight of facets that an earlier stage left
uncovered.  The first sum is the objective: its term for $\sigma$
vanishes exactly when some chain already covers $\sigma$, so the sum
vanishes exactly when every facet is covered, and $\beta$ grades the rest by how
far the \emph{best} chain is from covering them, which is the graded
signal that binary extension feedback lacks.  The second sum penalises defects in \emph{all}
chains, not only the best one; it is the cleanliness term that
suppresses drift (Section~\ref{sec:fail})---in the sphere runs
tested, $\gamma=0$ caused long chains to fail, and $\gamma=0.1$ was
insufficient.  (The later surface campaigns of Section~\ref{sec:descent}
use $\gamma=0.1$ with the additional stretch moves described there; we
did not retest this sentence for them.)  \emph{Coverage}, the quantity
reported in the tables, is the number of facets with $m(\sigma)=0$; it
is not itself a term of $E$.
Full coverage is equivalent to $m(\sigma)=0$ for every maximal
product simplex $\sigma$, not to vanishing total energy. The
all-chain defect penalty may remain positive at a complete cover.
The stopping criterion is therefore full coverage followed by
verification, never $E=0$.

The word \emph{energy} is the standard simulated-annealing term
\cite{KGV83,Cerny85}, not a physical quantity: in the Metropolis
dynamics \cite{Metropolis53}, the cost function of a combinatorial
search plays the role that energy plays in statistical mechanics, and
the temperature $T>0$ entering the Metropolis-style acceptance rule
$\min\{1,e^{-\Delta E/T}\}$ below is a control parameter measured in the same (dimensionless)
units as $E$.

A population of $S$ states anneals in parallel on the GPU under a
Metropolis-style acceptance rule of its own, distinct from the \textsc{LocalSearch}
kernel of the first phase (Section~\ref{sec:intro}).  The population is
not drawn with independent uniform interior maps: each state is built
by sweeps that flip entries of $\pi_1$ towards $\pi_2$ in random order,
so that every initial chain already interpolates between the two
projections.  A proposal chooses one chain and one vertex uniformly at random and
changes the value assigned to that vertex, either at a single interior
position or over a block of two to eight consecutive positions; the
block proposals are what we call \emph{track moves}, and they are drawn
with probability $p_{\mathrm{track}}$.  The new value is copied from a
map adjacent to the touched block with probability $1/2$, and otherwise
drawn uniformly from the remaining vertices; copy moves let a locally
good assignment propagate along the chain.  Copy proposals depend on
the current state, and a heterogeneous block can collapse to a constant
value that no single proposal of the same type reverts; symmetry of the
proposal kernel has not been established.  The acceptance rule alone
therefore does not give detailed balance with respect to a Boltzmann
distribution, and population selection and guided weights add further
adaptation.  We claim no equilibrium-sampling or global-convergence
guarantee for this heuristic.  Track moves are a move
shape, not an explicit bookkeeping of the facets a chain has been
assigned.  The temperature follows a geometric schedule \emph{within}
each epoch, and the incumbent best state is kept by highest unweighted
coverage, which need not be the state of lowest energy.  Every second
epoch the population undergoes selection and replacement---the weakest
quarter of the states is overwritten by copies of the elite and a few
slots are re-initialised at random---so the $S$ trajectories are not
independent for the whole run.  Once the incumbent covers every facet
the run stops, and the incumbent's chains are used to \emph{partition}
the facets (each facet goes to the first chain covering it).
This partitions the maximal facets. Each planner domain is the subcomplex
generated by its assigned facets; different domains may share
lower-dimensional faces.  A chain that receives no facet yields an empty
domain, which is discarded: $k$ chains produce at most $k$ nonempty
domains, and the bound $\SCst(K)\le k-1$ remains valid.  Each
restricted chain is then shortened by the greedy farthest-jump reduction of
\cite[Algorithm 2]{OLGB21}, and the resulting system is verified map by
map, link by link, facet by facet.  The output is a certificate, not a
claim (Section~\ref{sec:repro}).
Algorithm~\ref{alg:anneal} gives the full loop in the style of
\cite{OLGB21}.

\begin{algorithm}[H]
\caption{\textsc{ChainAnneal}($K$, $k$, $c_{\max}$; $S$, $\beta$,
$\gamma$, $p_{\mathrm{track}}$): anneal whole systems of $k$ chains over $K\times K$.}
\label{alg:anneal}
\begin{algorithmic}[1]
\Require a \emph{state} is $(C_1,\dots,C_k)$ with
$C_i=(\varphi_{i,0},\dots,\varphi_{i,c_{\max}})$,
$\varphi_{i,0}=\pi_1$, $\varphi_{i,c_{\max}}=\pi_2$, interior entries
arbitrary vertex assignments; population size $S$
\State initialise $S$ states by random-order sweeps that flip entries
of $\pi_1$ towards $\pi_2$
\State $\mathcal{C}^*\gets$ incumbent (highest coverage seen so far)
\For{epoch $=1,2,\dots$}
\ForAll{states $\mathcal{C}$ of the population, in parallel on the GPU}
\State perform $M$ annealing steps, the temperature $T$ following a
geometric schedule $T_0\to T_1$ \emph{inside} the epoch: pick one chain
$j$ and one vertex $w$ uniformly and propose a new value $u$ for
$\varphi_{j,i}(w)$, at a single interior position $i$ or (with
probability $p_{\mathrm{track}}$) over a block of two to eight
consecutive positions $i_1,\dots,i_2$; with probability $1/2$ the value
$u$ is copied from a neighbouring map, $\varphi_{j,i_1-1}(w)$ or
$\varphi_{j,i_2+1}(w)$, otherwise drawn uniformly from the remaining
vertices; accept with probability $\min\{1,e^{-\Delta E/T}\}$
\Statex \hspace{2em} with
$E=\sum_\sigma w_\sigma\,\mathrm{pen}(m(\sigma))
+\gamma\sum_j\sum_\sigma b_j(\sigma)$
\EndFor
\State update $\mathcal{C}^*$ by highest unweighted coverage
\Comment{not necessarily lowest $E$}
\If{$\mathcal{C}^*$ covers every facet of $K\times K$}
\State partition: each facet goes to the first chain of $\mathcal{C}^*$
covering it; discard empty domains \Comment{facet $\mapsto$ domain}
\State shorten each restricted chain by the greedy farthest-jump
reduction \cite[Algorithm~2]{OLGB21}
\State verify: every map simplicial and every link $1$-contiguous on
every facet of its domain, chain from $\pi_1|_J$ to $\pi_2|_J$
\State \Return $(\text{partition},\text{chains})$
\Comment{a certificate for $\SCst(K)\le k-1$}
\EndIf
\State every second epoch, overwrite the weakest quarter of the
population with copies of the elite and re-initialise a few slots at
random \Comment{selection and replacement}
\EndFor
\end{algorithmic}
\end{algorithm}

Tables~\ref{tab:annealer} and~\ref{tab:plateau} summarize the runs and Table~\ref{tab:params}
lists the search parameters of the complete runs.  The plateau runs are as
instructive as the successes, and we return to them in
Section~\ref{sec:products}: the two-chain runs on $\partial\Delta^3$,
$\partial\Delta^4$ and $\partial\Delta^5$ stop at coverage values that
are exactly those attained by suitable two-rectangle constructions.

\begin{table}[H]
\centering\footnotesize
\setlength{\tabcolsep}{3.5pt}
\begin{tabular}{@{}llr@{\hspace{8pt}}crl>{\raggedright\arraybackslash}p{3.0cm}@{}}
\toprule
target & triangulation $K$ & facets & chains $k$ & best coverage & wall time & outcome (domains found) \\
\midrule
$S^1$ & $\partial\Delta^2$ & $18$ & $2$ & $18/18$ & $<15$\,s
  & $9+9$, verified; $\SCst{\le}1$, so $\SCst{=}1{=}\TC$\\
$S^2$ & $\partial\Delta^3$ & $96$ & $3$ & $96/96$ (epoch $1$) & $22$\,s
  & $54+29+13$; $\SCst{=}2{=}\TC$\\
$S^3$ & $\partial\Delta^4$ & $500$ & $3$ & $500/500$ (epoch $1$) & $153$\,s
  & $309+130+61$; $\SCst{\le}2$\\
$S^4$ & $\partial\Delta^5$ & $2520$ & $3$ & $2520/2520$ (epoch $1$) & $13$\,min
  & $1634+653+233$; $\SCst{=}2{=}\TC$\\
$S^5$ & $\partial\Delta^6$ & $12348$ & $3$ & $12348/12348$ (epoch $2$) & $34$\,min
  & $7061+4239+1048$; $\SCst{\le}2$\\
$S^6$ & $\partial\Delta^7$ & $59136$ & $3$ & $59136/59136$ (epoch $2$) & ${\approx}2.5$\,h$^\S$
  & $44909+11451+2776$; $\SCst{=}2{=}\TC$\\
\midrule
$S^1\vee S^1$ & $S^1\vee S^1$ & $50$ & $3$ & $50/50$ (epoch $25$) & $41$\,s
  & $18+16+16$; $\SCst{=}2$\\
$S^2$ & cube boundary & $864$ & $3$ & $864/864$ (epoch $22$) & ${\approx}24$\,min$^\dagger$
  & $482+236+146$; $\SCst{=}2{=}\TC$\\
M\"obius band & $K_M$ & $150$ & $4$ & $150/150$ (epoch $1$) & $17$\,s$^\ddagger$
  & $54+37+34+25$; $\SCst{\le}3$\\
\bottomrule
\end{tabular}
\caption{Chain-first annealer, complete covers, one Tesla T4 unless noted.  The column $k$ is the
number of chains the search anneals; $k$ chains export at most $k$
nonempty domains, and $k$ is not a proven minimal number of domains.
On the simplex boundaries the star-rectangle constant is $\rho=3$
(Corollary~\ref{cor:spheres}); on the cube boundary $\rho=8$
(exact set cover, Table~\ref{tab:rho}), so the
three domains found use five fewer domains than any cover by the star
rectangles $D(v,w)$ of Theorem~\ref{thm:prod}.  The circle time is the
historical recorded value (Table~\ref{tab:params}), not a new
measurement.  Domain sizes are those of the
certified cover.  The per-run values of $c_{\max}$, the population size
and the move budget are \emph{not} uniform across rows; they are listed
in Table~\ref{tab:params}.  $^\dagger$The cube
entry is a three-rung ladder totalling about $24$ minutes---cold
rung ($303$\,s, $662/864$), warm rung ($814$\,s, $820/864$), guided
warm rung ($304$\,s, complete); rungs, warm starts and guided weights
are defined below.  The $304$\,s quoted above is the final guided rung
alone, not the whole sequence.  $^\ddagger$Warm-started from a certified
$5$-domain cover ($k{=}5$ cold, $36$\,s, epoch~$3$).  Eight independent
$k{=}3$ searches (cold, warm-restarted and guided) all stalled at
$148/150$ facets; a companion note \cite{PM26} reports a $3$-domain
cover found by constraint programming.
$^\S$Eight-rung ladder on an NVIDIA RTX PRO 6000 Blackwell Server Edition
(not the T4 used for all other rows): cold $5$-rung phase
(${\approx}1$\,h\,$42$\,min, $S{=}64$), then guided $3$-rung
continuation (${\approx}47$\,min, $S{=}256$).}
\label{tab:annealer}
\end{table}

\begin{table}[H]
\centering\footnotesize
\setlength{\tabcolsep}{3.5pt}
\begin{tabular}{@{}llr@{\hspace{8pt}}crl>{\raggedright\arraybackslash}p{3.0cm}@{}}
\toprule
target & triangulation $K$ & facets & chains $k$ & best coverage & wall time & outcome \\
\midrule
$S^2$ & $\partial\Delta^3$ & $96$ & $2$ & $84/96$ & (plateau)
  & facet coverage; matches the two-star-rectangle benchmark\\
$S^3$ & $\partial\Delta^4$ & $500$ & $2$ & $460/500$ & (plateau)
  & facet coverage; matches the two-star-rectangle benchmark\\
$S^4$ & $\partial\Delta^5$ & $2520$ & $2$ & $2380/2520$ (epoch $1$) & $12$\,min
  & facet coverage; matches the two-star-rectangle benchmark\\
$S^2$ & cube boundary & $864$ & $2$ & --- & ---
  & $\TC(S^2)=2$: no $2$-domain cover exists\\
\bottomrule
\end{tabular}
\caption{Two-chain plateau runs ($k=2$): the search seeks a cover by
$2$ domains and stalls.  Coverage counts maximal simplices of the
ordered product. Equality with a two-star-rectangle coverage fraction
does not imply that the same fraction of base-facet pairs is covered in
its entirety; for example, \texttt{s3\_k2\_v2\_point} covers $21$ of
$25$ base-facet pairs completely and \texttt{s4\_k2\_v2} covers $33$ of
$36$.  The plateau values recurred across the seeds and continuations
we ran; these are observations on a finite set of runs, not a
probability estimate.}
\label{tab:plateau}
\end{table}

\begin{table}[H]
\centering\footnotesize
\begin{tabular}{@{}lrrrr@{}}
\toprule
run & $c_{\max}$ & $S$ & moves/epoch & seed \\
\midrule
$\partial\Delta^3$ (S$^2$) & $20$ & $4096$ & $20000$ & $51$\\
$\partial\Delta^4$ (S$^3$) & $40$ & $4096$ & $20000$ & $52$\\
$\partial\Delta^5$ (S$^4$) & $40$ & $4096$ & $20000$ & $61$\\
$\partial\Delta^6$ (S$^5$) & $16$ & $4096$ & $5000$ & $71$\\
$\partial\Delta^7$ (S$^6$, final guided rung) & $56$ & $256$ & --- & $608$\\
wedge $S^1\vee S^1$ & $16$ & $4096$ & $5000$ & $71$\\
cube boundary (final guided stage) & $16$ & $4096$ & $5000$ & $84$\\
M\"obius band ($k{=}4$, warm) & $40$ & $4096$ & $5000$ & $4001$\\
\bottomrule
\end{tabular}
\caption{Search parameters of the runs reported in
Table~\ref{tab:annealer}; the first seven rows used $\beta=5$, $\gamma=0.5$ and
$p_{\mathrm{track}}=0.3$.  The $S^6$ entry is the final rung of the
$8$-rung ladder described in Table~\ref{tab:annealer}; earlier rungs
escalated $c_{\max}$ from $16$ to $56$ and used $S{=}64$ in the cold
phase ($750$ iterations~$\times$ $4$ epochs, full coverage at epoch~$2$).  These settings should not be read back onto
the circle run, which used $k=2$, $c_{\max}=40$, $S=1024$, $20000$
moves per epoch, seed $5$, $\beta=5$, $\gamma=0.5$ and
$p_{\mathrm{track}}=0.05$ (recorded wall time $7.9$\,s; certificate
\texttt{chain-first-annealer/opti/results/circle\_probe3}, chains of
$4$ and $5$ links), nor onto the M\"obius band run, which used
$\beta=1$, $\gamma=0.5$ and $p_{\mathrm{track}}=0.5$.}
\label{tab:params}
\end{table}

\emph{Ladders of restarts: rungs, warm starts, guided weights.}  The
runs of Table~\ref{tab:annealer} were single calls to
Algorithm~\ref{alg:anneal}, with one exception.  The cube boundary
certificate was reached by a short \emph{ladder}: a sequence of bounded
calls, each of which we call a \emph{rung}, chained so that each rung
starts where the previous one stopped.  A \emph{cold} rung is an
ordinary call to Algorithm~\ref{alg:anneal}; on the cube it covered
$662$ of $864$ facets in $303$\,s.  A \emph{warm} rung initialises its
population from the previous rung's incumbent: half of the population
receives the $k$ incumbent chains, each shortened on its domain by the
greedy reduction and padded with copies of $\pi_2$ up to $c_{\max}$
links, and the other half is built from fresh sweeps as before; on the
cube it continued to $820/864$ in $814$\,s.  A \emph{guided} rung
additionally rebalances the energy during the search: after each epoch
every weight decays back towards $1$ by $w_\sigma\gets
1+0.7\,(w_\sigma-1)$, and every facet the incumbent still misses
receives an additive increment $g$, capped at $g_{\max}$; reported
coverage stays unweighted.  A guided warm rung closed the remaining
$44$ facets in $304$\,s.  In the configurations tested, guidance without the seed was not
sufficient: a guided \emph{cold} rung stalled at $858/864$ with the
whole population pinned at $856$--$858$---the guide needed a good
incumbent, not only a budget.

Why short rungs at all?  The coverage gained within a single warm rung
is \emph{front-loaded}: an empirical fit to the dodecahedron campaign
below gives $\Delta\mathrm{cov}(t)\approx c\,t^{0.8}$ (one
rung gained $+33$ facets in its first $7$ minutes, $+58$ in the first
$15$, and $+170$ in the full hour).  Early in a rung the population
claims the facets accessible from the incumbent's chain topology; once
those are exhausted, further time buys rapidly diminishing returns.
Every rung also pays a fixed overhead $\tau$ (kernel compilation and
host--device transfer, $\tau\approx40$\,s on the machine used here), so
the coverage harvested per GPU-hour, $c\,t^{0.8}/(t+\tau)$, peaks at
$T^{*}_{\mathrm{rung}}=\tfrac{0.8}{1-0.8}\,\tau=4\tau\approx2.7$\,min.
Algorithm~\ref{alg:ladder} abstracts the pattern.  It is a ladder at
fixed $k$; the surface campaigns also used a ladder in $k$ itself,
descending one unit at a time from a complete cover (Section~\ref{sec:descent}).

\begin{algorithm}[H]
\caption{\textsc{WarmLadder}($K$, $k$, $c_{\max}$; $S$, $\beta$,
$\gamma$, $p_{\mathrm{track}}$; per-rung budget $T_{\mathrm{rung}}$,
guided increment $g$, cap $g_{\max}$, max rungs $R_{\max}$).
Algorithm~\ref{alg:anneal} run as a ladder: warm seeding between rungs,
guided reweighting within each rung.}
\label{alg:ladder}
\begin{algorithmic}[1]
\State run \textsc{ChainAnneal}$(K,k,c_{\max};S,\beta,\gamma,
p_{\mathrm{track}})$ for at most $T_{\mathrm{rung}}$ (\emph{cold} rung,
$w_\sigma\equiv1$); set $\mathcal{C}^{*}\gets$ best state found
\For{$r=1,2,\ldots,R_{\max}$}
  \If{$\mathcal{C}^{*}$ covers all facets of $K\times K$}
    \State\Return $\mathcal{C}^{*}$
    \Comment{certificate: $\SCst(K)\le k-1$}
  \EndIf
  \State \emph{warm seeding}: half the population $\gets$ the chains of
  $\mathcal{C}^{*}$, each reduced on its domain and padded with $\pi_2$
  up to $c_{\max}$ links; the other half is freshly swept
  \State \emph{guided rung}: run \textsc{ChainAnneal} for at most
  $T_{\mathrm{rung}}$ with adaptive weights---after each epoch,
  $w_\sigma\gets1+0.7\,(w_\sigma-1)$ for every facet $\sigma$ and
  $w_\sigma\gets\min(w_\sigma+g,\,g_{\max})$ for every facet the
  incumbent does not cover; update $\mathcal{C}^{*}\gets$ best state
  found \Comment{coverage reporting stays unweighted}
\EndFor
\State\Return $\mathcal{C}^{*}$
\Comment{partial cover; no certificate}
\end{algorithmic}
\end{algorithm}

To probe how far ladders extend the reach of the search, we applied
Algorithm~\ref{alg:ladder} to the dodecahedron boundary: a fan
triangulation of $S^2$ with $20$ vertices and $36$ triangles, yielding
$K\times K$ with $7776$ product facets.  Its simplex density is $3.2\%$
(compared to $21.4\%$ for the cube boundary), and the dual graph of $K$
has diameter $7$ against $4$ for the cube---both figures suggest a
substantially harder search.  We ran five long warm escalations of up
to one hour each, then a systematic $2.5$-minute ladder of $172$ short
rungs executed unattended over $12$ GPU-hours.
Table~\ref{tab:dodeca} records the complete campaign.

\begin{table}[H]
\centering\footnotesize
\begin{tabular}{@{}lrrrr@{}}
\toprule
stage & time & coverage & gain & \% \\
\midrule
cold start ($c_{\max}=40$, $S=4096$) &
  $60$\,min & $0\to 644$ & $+644$ & $8.3$ \\
warm escalation 1 &
  $60$\,min & $644\to 1208$ & $+564$ & $15.5$ \\
warm escalation 2 &
  $60$\,min & $1208\to 1579$ & $+371$ & $20.3$ \\
warm escalation 3 &
  $60$\,min & $1579\to 1749$ & $+170$ & $22.5$ \\
warm escalation 4 &
  $15$\,min & $1749\to 1802$ & $+53$ & $23.2$ \\
warm escalation 5 &
  $7$\,min & $1802\to 1831$ & $+29$ & $23.5$ \\
\midrule
2.5-min ladder, R.\,1--20 &
  $55$\,min & $1831\to 2064$ & $+233$ & $26.5$ \\
2.5-min ladder, R.\,22--41 &
  $58$\,min & $2075\to 2218$ & $+143$ & $28.5$ \\
2.5-min ladder, R.\,42--61 &
  $58$\,min & $2218\to 2310$ & $+92$ & $29.7$ \\
2.5-min ladder, R.\,62--81 &
  $58$\,min & $2310\to 2400$ & $+90$ & $30.9$ \\
2.5-min ladder, R.\,82--101 &
  $58$\,min & $2400\to 2441$ & $+41$ & $31.4$ \\
2.5-min ladder, R.\,102--121 &
  $58$\,min & $2441\to 2487$ & $+46$ & $32.0$ \\
2.5-min ladder, R.\,122--141 &
  $58$\,min & $2487\to 2523$ & $+36$ & $32.4$ \\
2.5-min ladder, R.\,142--172 + probes &
  ${\approx}3.5$\,h & $2523\to 2602$ & $+79$ & $33.5$ \\
\midrule
\textbf{total} &
  ${\approx}20$\,h & $0\to 2602$ & $+2602$ & $33.5$ \\
\bottomrule
\end{tabular}
\caption{Complete warm-start rung-ladder campaign on the dodecahedron boundary
($7776$ product facets of $K\times K$, $k=3$).
Upper block: five long warm escalations ($c_{\max}=40$, $g=3$, $g_{\max}=11$),
each reseeded from the previous incumbent.
Lower block: $2.5$-minute ladder run unattended over $12$\,GPU-hours,
summarised in batches of $20$ rungs; the final row combines the last partial
batch with extended window probes and a $70$-trial basin survey.
Mean gain per rung falls from $11.7$ in the first batch to below $1$ in
the last; no complete cover was found.
We conjecture that the obstacle is the sparsity of $K$: a dual-graph
diameter of $7$ may require coordinated walks that become harder to
satisfy simultaneously as coverage rises.  This is a hypothesis; an
unsuccessful campaign does not establish the cause.}
\label{tab:dodeca}
\end{table}

We close the section with four observations on
Tables~\ref{tab:annealer} and~\ref{tab:plateau}.

\emph{Certified values, not just bounds.}  The successes pin the
invariant itself: for $S^2$ and $S^4$ the covers found, together with
the lower bound $\TC=2$ of Section~\ref{sec:products}, force
$\SCst=\TC=2$, so the search independently reproduces at level $0$ the
optimum that theorem proves in general.

\emph{Near-product shape.}  The winning domains on the simplex
boundaries are close in size to star rectangles but not contained in
them.  On $\partial\Delta^4$ the certified partition ($309+130+61$) has
its largest domain at $309$ facets where a single rectangle holds $320$;
on $\partial\Delta^5$ the certified partition ($1634+653+233$) has its
largest at $1634$ where a rectangle holds $1750$.  The domains overlap
with rectangles but are not unions of rectangles.

\emph{The plateau value is structural; the plateau shape is not.}
The stalled two-chain runs stop exactly at the coverage of the best
two-rectangle constructions.  This does not mean the search found
rectangles: among the saved states, one $\partial\Delta^4$ run ends in
two genuine rectangles of $320$ facets each, union $460$, while two
other runs reach the same $460$ with domains contained in no rectangle
$\operatorname{star}(v)\times\operatorname{star}(w)$ at all, and the
$\partial\Delta^5$ run at $2380/2520$ is likewise not a union of two
rectangles.\footnote{The runs cited are \texttt{s3\_k2\_v2\_track}
(two rectangles of $320$), \texttt{s3\_k2\_A} (domains of $284$ and
$265$, the second in no rectangle) and \texttt{s3\_k2\_v2\_point}
(domains of $297$ and $317$, neither in a rectangle; union still
$460$) on $\partial\Delta^4$, and \texttt{s4\_k2\_v2} ($2380/2520$)
on $\partial\Delta^5$.  These saved states are partial states, not
certificates; they are in \texttt{chain-first-annealer/opti/results/}
(Section~\ref{sec:repro}).}  What is stable across seeds is the coverage
value, not the domains that realise it---and explaining \emph{that
number} is what makes the next section possible.

\emph{Non-product freedom pays.}  The cube boundary, an $8$-vertex,
$12$-face triangulation of $S^2$ distinct from $\partial\Delta^3$,
shows the extreme case: the star-rectangle bound $\rho=8$ certifies
only $\SCst\le7$, while the free-form three-chain search reaches
$\SCst=2=\TC(S^2)$: eight star-rectangle domains versus three
general contiguity domains, five domains fewer than any cover by the
star rectangles $D(v,w)$ of Theorem~\ref{thm:prod}.

\section{The product domains theorem and the star cover}\label{sec:products}

Write $\operatorname{star}(v)$ for the family of facets of $K$ containing
the vertex $v$ (not the open star), and $c_x$ for the constant map to the vertex $x$.

\begin{theo}[product domains]\label{thm:prod}
Let $K$ be a finite, nonempty, connected simplicial complex with
ordered vertex set.  For vertices $v,w$ of $K$, let $D(v,w)\subseteq K\times K$ be the
subcomplex generated by the maximal facets lying over pairs $(s,t)$ of
facets of $K$ with $v\in s$ and $w\in t$.  Then $D(v,w)$ is a
contiguity domain for
$(\pi_1,\pi_2)$: if $v=v_0,v_1,\dots,v_m=w$ is a path in the
$1$-skeleton of $K$ (such a path exists since $K$ is connected), then on $D(v,w)$
\[
\pi_1 \;\sim\; c_{v_0}\;\sim\; c_{v_1}\;\sim\;\cdots\;\sim\;
c_{v_m}\;\sim\; \pi_2,
\]
a chain of $m+2$ links ($2$ links when $v=w$).
\end{theo}

\begin{proof}
The generating maximal facets are maximal in $K\times K$, so none is a
face of another; the facets of the subcomplex they generate are exactly
the generators, and since $1$-contiguity is a condition on facets
(Section~\ref{sec:back}), it suffices to check the chain on those.  A
facet $\sigma\in D(v,w)$ lying over $(s,t)$ has $\pi_1(\sigma)=s$,
$\pi_2(\sigma)=t$ (a staircase path visits every row and every column,
so all vertices of $s$ and of $t$ occur), and $c_x(\sigma)=\{x\}$, so
every map in the chain is simplicial on $\sigma$.  Consecutive maps are
$1$-contiguous on $\sigma$: $\pi_1(\sigma)\cup c_v(\sigma)=s\cup\{v\}=s$
since $v\in s$; $c_{v_i}(\sigma)\cup c_{v_{i+1}}(\sigma)=\{v_i,v_{i+1}\}$
is an edge of $K$; and $c_w(\sigma)\cup\pi_2(\sigma)=t$.
\end{proof}

We emphasize the provenance: this statement was not guessed, it was
\emph{found}.  The plateaus of Table~\ref{tab:plateau} sit at exactly
the coverage two such rectangles achieve---though, as recorded in
Section~\ref{sec:annealer}, the domains realising that coverage need not
themselves be rectangles---and in naming the value the theorem converts
a search problem into a covering problem.

\begin{corollary}[a search-free upper bound]\label{cor:rho}
Let $\rho(K)$ be the minimum number of rectangles $D(v,w)$ covering all
facet pairs of $K\times K$, i.e.\ the minimum size of a set
$R\subseteq V\times V$ such that for every pair $(s,t)$ of facets of
$K$ there is $(v,w)\in R$ with $v\in s$, $w\in t$.  Then
\[
\SCst(K)\;\le\;\rho(K)-1.
\]
\end{corollary}

\begin{proof}
Every facet pair $(s,t)$ has some $(v,w)\in R$ with $v\in s$, $w\in
t$, so the domains $D(v,w)$, $(v,w)\in R$, cover $K\times K$ (covers
need not be disjoint); Theorem~\ref{thm:prod} supplies the chains, and
$\rho(K)$ domains certify $\SCst(K)\le\rho(K)-1$.
\end{proof}

The number $\rho(K)$ is the optimum of a set-cover instance with
$|V|^2$ candidate sets over $|F|^2$ elements, and it is computed
exactly, not heuristically.  The repository's \texttt{rect\_cover.py}
does so by bitmask branch and bound; for this paper we re-derived every
value of Table~\ref{tab:rho} with the CP-SAT solver of OR-Tools
\cite{ORTools}, which returned status OPTIMAL (a feasible cover and a
matching lower bound) on all fifteen instances (the twelve rows of the table, with the simplex
boundaries $n=1,\dots,4$ computed separately), in at most $12$\,s on four
CPU threads.  Exactness matters: a greedy choice of rectangles needs
$27$ rectangles on the icosahedron boundary against $\rho=24$, $13$ on
the Cs\'asz\'ar torus against $11$, and $16$ on the hemi-dodecahedron
against $15$.

\begin{corollary}[all boundary-of-simplex spheres]\label{cor:spheres}
For every $n\ge1$, $\rho(\partial\Delta^{n+1})=3$; hence
\[
\SCst(\partial\Delta^{n+1})\le 2\ \ \text{for all }n\ge1,
\quad
\SCst(\partial\Delta^{n+1})=\TC(S^n)=2\ \ \text{for even }n\ge2,
\]
sharp at subdivision level $0$.
\end{corollary}

\begin{proof}
The facets of $\partial\Delta^{n+1}$ are $s_a=V\setminus\{a\}$,
$a\in\{0,\dots,n+1\}$.  Take $R=\{(0,0),(1,1),(2,2)\}$: the pair
$(s_a,s_b)$ is missed by $(v,v)$ iff $v\in\{a,b\}$, so it is missed by
all three rectangles iff $\{0,1,2\}\subseteq\{a,b\}$, impossible.  Two
rectangles $(v,w),(v',w')$ never suffice: $(s_v,s_{w'})$ is missed by
both.  Hence $\rho=3$, and for even $n$ Farber's $\TC(S^n)=2$
\cite{F03} (reduced convention) with $\TC(\|K\|)\le\SCst(K)$
\cite{G18} gives equality.  The argument excludes $n=0$: $S^0$ is
disconnected.
\end{proof}

\begin{lemma}[counting bounds]\label{lem:rows}
Let $\tau$ be the minimum size of a set of vertices meeting every facet
of $K$ (a \emph{transversal}), let $d=\max_v|\operatorname{star}(v)|$,
and let $F$ be the number of facets of $K$.  Then
\[
\Bigl\lceil \frac{\tau F}{d}\Bigr\rceil\;\le\;\rho(K)\;\le\;\tau^{2}.
\]
\end{lemma}

\begin{proof}
\emph{Upper bound.}  Let $S$ be a transversal with $|S|=\tau$.  For a
pair $(s,t)$ of facets choose $v\in s\cap S$ and $w\in t\cap S$; then
$(s,t)$ lies over $(v,w)\in S\times S$, so the $\tau^2$ rectangles
$D(v,w)$, $v,w\in S$, cover $K\times K$.

\emph{Lower bound.}  Fix a set $R$ of $\rho(K)$ rectangles covering all
facet pairs and a facet $s$, and put
$W_s=\{w:\exists\,v\in s\text{ with }(v,w)\in R\}$.  Every facet $t$ meets
$W_s$ (the pair $(s,t)$ lies over some $(v,w)\in R$ with $v\in s$ and
$w\in t$), so $W_s$ is a transversal and $|W_s|\ge\tau$.  Counting the
pairs $(s,(v,w))$ with $(v,w)\in R$ and $v\in s$ in two ways,
\[
\sum_{(v,w)\in R}|\operatorname{star}(v)|
=\sum_s\#\{(v,w)\in R:v\in s\}
\;\ge\;\sum_s|W_s|\;\ge\;\tau F ,
\]
where the first inequality can be strict, because two vertices of $s$
may be paired with the same $w$.  Every summand on the left is at
most $d$, so $d\,\rho(K)\ge\tau F$.
\end{proof}

Lemma~\ref{lem:rows} is elementary, but it is not vacuous: it computes
$\rho$ outright in several cases, and it is the only general
tool for $\rho$ that we have beyond computation.  It gives $\rho(\partial\Delta^{n+1})\ge3$
for every $n\ge1$ (here $\tau=2$, $d=n+1$, $F=n+2$, so
$\lceil2(n+2)/(n+1)\rceil=3$), the exact value $\rho=6$ of the
hemi-icosahedron ($\tau=3$, $d=5$, $F=10$), and, most strikingly, the
exact value $\rho=24$ of the icosahedron boundary: $\tau=6$, $d=5$,
$F=20$ give $\lceil 6\cdot20/5\rceil=24$, matched by an explicit
$24$-rectangle cover.  Table~\ref{tab:rho} lists both bounds next to the
exact $\rho$ for the complexes of our campaigns; the sandwich holds in every row, the
lower bound is attained in five rows and the upper bound in four.

\begin{table}[H]
\centering\footnotesize
\setlength{\tabcolsep}{4pt}
\begin{tabular}{@{}lrrrrrrr@{}}
\toprule
complex $K$ & $|V|$ & $F$ & $\tau$ & $d$ &
$\lceil\tau F/d\rceil$ & $\rho(K)$ & $\tau^2$ \\
\midrule
$\partial\Delta^{n+1}$, $n\ge1$ (incl.\ the circle) & $n{+}2$ & $n{+}2$ & $2$ & $n{+}1$ & $3$ & $3$ & $4$\\
wedge $S^1\vee S^1$ & $4$ & $5$ & $2$ & $3$ & $4$ & $4$ & $4$\\
M\"obius band $K_M$ & $5$ & $5$ & $2$ & $3$ & $4$ & $4$ & $4$\\
cube boundary & $8$ & $12$ & $3$ & $6$ & $6$ & $8$ & $9$\\
hemi-icosahedron ($\mathbb{RP}^2$) & $6$ & $10$ & $3$ & $5$ & $6$ & $6$ & $9$\\
hemicube ($\mathbb{RP}^2$) & $7$ & $12$ & $3$ & $6$ & $6$ & $7$ & $9$\\
hemi-octahedron ($\mathbb{RP}^2$) & $9$ & $16$ & $3$ & $6$ & $8$ & $9$ & $9$\\
hemi-dodecahedron ($\mathbb{RP}^2$) & $10$ & $18$ & $4$ & $9$ & $8$ & $15$ & $16$\\
Cs\'asz\'ar torus & $7$ & $14$ & $4$ & $6$ & $10$ & $11$ & $16$\\
Klein bottle (Cervone) & $8$ & $16$ & $4$ & $7$ & $10$ & $12$ & $16$\\
icosahedron boundary & $12$ & $20$ & $6$ & $5$ & $24$ & $24$ & $36$\\
dodecahedron boundary & $20$ & $36$ & $6$ & $8$ & $27$ & $36$ & $36$\\
\bottomrule
\end{tabular}
\caption{The rectangle constant $\rho(K)$ and the bounds of
Lemma~\ref{lem:rows}.  Here $\tau$ is the transversal number, $d$ the
largest vertex star and $F$ the number of facets; $\rho$ is the exact
optimum (CP-SAT, status OPTIMAL, recomputed from the facet lists stored
in \texttt{results/*\_starcover*/summary.json}).  The triangulations are
the ones of the repository; the hemicube and the hemi-octahedron are
subdivided to be simplicial (a centre vertex per square face, a
midpoint per edge cell), which is why $|V|=7$ and $9$, and the pentagons
of the hemi-dodecahedron are fanned without new vertices.  The value $\rho=3$ for
$\partial\Delta^{n+1}$ is Corollary~\ref{cor:spheres} and was also
recomputed for $n=1,\dots,4$.  The coincidence $\rho=|V|$ for the three
$\mathbb{RP}^2$ triangulations with $|V|=6,7,9$ is a property of these
triangulations (the hemi-dodecahedron has $\rho=15>10$) and has no
structural role.}
\label{tab:rho}
\end{table}

\subsection{The star cover as a method}\label{sec:starmethod}

Everything above is mathematics; what we overlooked for a long time is
that it is also the best first move.  By the \emph{star cover} of a
connected complex $K$ we mean the following procedure
(\texttt{src/star\_cover.py} and \texttt{src/rect\_cover.py}).

\begin{algorithm}[H]
\caption{\textsc{StarCover}($K$): a complete, verified system of
planners with no search.}
\label{alg:star}
\begin{algorithmic}[1]
\State list the facets of $K$; for every vertex pair $(v,w)$ form the
set of facet pairs $\operatorname{star}(v)\times\operatorname{star}(w)$
\State solve the set cover exactly: $R\gets$ a minimum set of vertex
pairs whose rectangles contain all $|F|^2$ facet pairs
\Comment{$|R|=\rho(K)$}
\ForAll{$(v,w)\in R$}
  \State chain $\pi_1\sim c_{v_0}\sim\cdots\sim c_{v_m}\sim\pi_2$ along a
  shortest edge path $v=v_0,\dots,v_m=w$ of the $1$-skeleton
\EndFor
\State assign every facet of $K\times K$ to the first rectangle that
contains it; shorten each chain on its domain by the greedy reduction;
verify map by map, link by link, facet by facet
\State \Return the partition, the chains, and $\SCst(K)\le\rho(K)-1$;
store the chains (\texttt{best\_chains}) as a seed for
Algorithm~\ref{alg:anneal}
\end{algorithmic}
\end{algorithm}

The procedure cannot fail.  Completeness of the rectangle family is
immediate (a facet pair $(s,t)$ lies over $(v,w)$ for any $v\in s$,
$w\in t$), and the chains need no search because contiguity of every link
reduces to the membership statements $v\in s$, $w\in t$ in the proof of
Theorem~\ref{thm:prod}.  Its cost is that of a small set cover plus a
linear pass over the facets of $K\times K$.  On a laptop CPU, a whole
run of the script \texttt{star\_cover.py} (greedy rectangle choice) on the
icosahedron boundary ($2400$ facets) took $1.3$\,s, and the exact set cover took $12$\,s
for the icosahedron and $4.6$\,s for the dodecahedron ($7776$ facets);
the simplex boundaries up to $\partial\Delta^9$ ($1\,287\,000$ facets)
are covered by three domains
(\texttt{results/s2\_starcover}\,\dots\,\texttt{s8\_starcover}), and all
of these certificates pass the independent auditor of
Section~\ref{sec:repro}.  We state the outputs, because each is used
later: (a) a canonical complete cover of $K\times K$; (b) an
unconditional, exact-in-family upper bound $\rho(K)-1$; (c) the chains
of the cover, which seed the annealer through its \texttt{--init}
option (Section~\ref{sec:descent}); and (d) a floor under the search:
$k\ge\rho(K)$ domains exist by explicit certificate, so the annealer
has something to find only below $\rho$.

\subsection{Why $\rho-1$ is only an upper bound, and when it is the
minimum}\label{sec:sharp}

The star cover minimises over a restricted family: domains of the form
$\operatorname{star}(v)\times\operatorname{star}(w)$, with chains of
constant maps, which factor through a point and are the most degenerate
chains possible.  General contiguity domains have arbitrary shapes and
general chains.  Hence $\SCst(K)\le\rho(K)-1$ and the whole gap lives in
the restriction to this family.  The gap is real, and large on sparse
complexes (Table~\ref{tab:surfaces}): the circle ($\rho-1=2$ against
$\SCst=1$), the wedge ($3$ against $2$), the M\"obius band ($3$ against
at most $2$), the cube boundary ($7$ against $2$) and the icosahedron
boundary ($23$ against at most $4$).

The bound is the minimum exactly where a lower bound rises to meet it.
This happens on the even spheres: for every even $n\ge2$,
\[
\TC(S^n)=2\;\le\;\SCst(\partial\Delta^{n+1})\;\le\;\rho-1=2,
\]
which is Corollary~\ref{cor:spheres}; the star cover is minimal there
not because the product family is exact in general but because the
topological floor happens to sit at the ceiling of the crudest
constructive family.  The two neighbouring cases show that this is a
coincidence of the complex and not a property of the method.  On the
circle, $\TC=1<2$ and the two-domain cover of Table~\ref{tab:annealer}
beats $\rho-1$.  On the odd spheres $n\ge3$ the sandwich does not close,
$\TC(S^n)=1<2=\rho-1$, and whether $\SCst(\partial\Delta^{n+1})$ is $1$
or $2$ is the subject of Section~\ref{sec:sat} and Conjecture~\ref{conj}.
The recorded annealed two-chain searches on $\partial\Delta^4$ did not
exceed within their budgets the coverage $460/500$ that two rectangles
already achieve, whatever the shape of the domains they returned.
Finally, ``sharp'' in Corollary~\ref{cor:spheres} refers to the
specific triangulation $\partial\Delta^{n+1}$ at subdivision level
$0$; in the barycentric limit the invariant is $\TC$ itself
\cite{G18}.

\section{Descents and seeds}\label{sec:descent}

The ladder of Section~\ref{sec:annealer} (Algorithm~\ref{alg:ladder})
holds $k$ fixed and restarts the search at the same $k$.  The surface
campaigns used a different pattern, which we call a \emph{$k$-descent
ladder}.  A complete $k$-domain cover is the warm start of the
$(k{-}1)$-domain search: the annealer keeps the $k{-}1$ largest domains
of the incumbent, each with its reduced chain padded with copies of
$\pi_2$ to $c_{\max}$ links, drops the smallest, and tries to absorb the
facets it covered.  The first rung is the expensive one, since it
needs a first cover.  The campaigns ran the annealer in the variant
\texttt{annealer\_stretch.py}: the energy of
Section~\ref{sec:annealer} unchanged, with $\beta=5$, $\gamma=0.1$,
$c_{\max}=40$, $S=4096$, $5000$ annealing steps per epoch, seed $8233$
in every descent below except the $\mathbb{RP}^2$ one, plus a compound ``stretch'' proposal (a block
of interior maps is set so that every vertex of one uncovered facet takes
its neighbour's value, with a ``collapse'' variant that carries a single
fresh value), a warm-start option \texttt{--init} that reads
\texttt{best\_chains} from a previous \texttt{summary.json}, and the
selection step of Algorithm~\ref{alg:anneal}.  Every wall time quoted
in this section is read from the \texttt{summary.json} or the log of the
run; the summaries do not attest the hardware, which was not uniform
across the campaigns.

\subsection{Annealer alone: start high, descend}

With the annealer alone, the working strategy was to start at a large
$k$, where a first complete cover is easiest to find, and to descend one
unit at a time.  The recorded descents are
$\mathbb{RP}^2$ ($k=11\to6$), the Klein bottle ($k=15\to10$) and the
hemicube ($k=20\to7$) in Table~\ref{tab:descent}.  The descent works:
once a cover exists, successive rungs are often cheap (several rungs
of the icosahedron ladder close within three epochs).  But the first,
cold, high-$k$ cover is not cheap, and it is the stage that does not scale:
it took two hours on the Klein bottle ($1536$ facets), $41$ minutes on
the hemicube ($864$ facets), and on the hemi-dodecahedron ($1944$ facets)
and hemi-octahedron ($1536$ facets) the cold runs we recorded at
$k=15,20,30$ and $k=30$ ended without a cover.  The latter four
are log files only (\texttt{logs/hemi\_dodeca\_k15.log}, \dots, \texttt{k30},
\texttt{logs/hemi\_octa\_k30.log}); they stop between $28$ and $188$
epochs with no certificate and with no recorded stopping criterion, so
they show that the cold start had not succeeded within the time spent,
$0.8$ to $4.1$ hours, not that it cannot.

\subsection{Star-seeded descent}

The star cover of Section~\ref{sec:starmethod} replaces that first
rung.  Algorithm~\ref{alg:star} supplies a complete $\rho$-domain cover
in seconds, and the annealer starts at $k=\rho-1$ from the $\rho-1$
largest star domains.  The recorded star-seeded runs are the hemi-octahedron
($k=8$, one rung, complete at epoch $106$), the hemi-dodecahedron
($k=14\to7$, every rung closed) and the icosahedron
($k=23\to5$, every rung closed; Section~\ref{sec:icosa}).  In the cleanest
comparison available, the hemi-dodecahedron, the star-seeded $k=14$ run
reached a complete cover in $6689$\,s, while cold runs at the larger
$k=15,20,30$ had not reached one after $3521$, $14\,673$ and $2863$\,s.
Both statements are recorded observations, single seed, one lineage per
complex; the comparison is not a controlled experiment (the cold runs
were interrupted rather than run to a fixed budget) and does not by
itself separate the effect of the seed from that of the continuation.
It is nevertheless the reason we now start every campaign with
Algorithm~\ref{alg:star}.

\begin{table}[H]
\centering\footnotesize
\setlength{\tabcolsep}{3pt}
\begin{tabular}{@{}l r>{\raggedright\arraybackslash}p{2.9cm}>{\raggedright\arraybackslash}p{2.8cm}>{\raggedright\arraybackslash}p{1.8cm}>{\raggedright\arraybackslash}p{2.1cm}@{}}
\toprule
complex ($|K{\times}K|$) & $\rho$ & first cover & descent (total wall) &
lowest verified & first stall \\
\midrule
$\mathbb{RP}^2$, hemi-icosa ($600$)
 & $6$ & cold $k{=}11$: $590/600$ in $623$\,s, then $56$\,s to complete
 & $k{=}11\to6$ ($0.8$\,h) & $k{=}6$, $\SCst\le5$
 & $k{=}5$: $540$--$544/600$\\
hemicube ($864$)
 & $7$ & cold $k{=}20$: $2470$\,s ($45$ epochs)
 & $k{=}20\to7$ ($2.0$\,h) & $k{=}7$, $\SCst\le6$
 & $k{=}6$: $828$, $836/864$\\
Klein bottle ($1536$)
 & $12$ & cold $k{=}15$: $1513/1536$ in $5675$\,s, then $1400$\,s to complete
 & $k{=}15\to10$ ($4.3$\,h) & $k{=}10$, $\SCst\le9$
 & $k{=}9$: $1494/1536$ (four chained runs)\\
hemi-octahedron ($1536$)
 & $9$ & cold $k{=}30$: $1335/1536$ after $7492$\,s, no cover (log)
 & star-seeded $k{=}8$: $6298$\,s & $k{=}8$, $\SCst\le7$
 & $k{=}7$: $1432/1536$, $13$ epochs (log)\\
hemi-dodecahedron ($1944$)
 & $15$ & cold $k{=}15,20,30$: $1324$, $1896$, $1345/1944$, no cover (logs)
 & star-seeded $k{=}14\to7$ ($12.3$\,h) & $k{=}7$, $\SCst\le6$
 & $k{=}6$: $1918$, $1922/1944$ ($300$ epochs each)\\
icosahedron ($2400$)
 & $24$ & star-seeded $k{=}23$: $1545$\,s, $23$ epochs
 & $k{=}23\to5$ ($17.9$\,h) & $k{=}5$, $\SCst\le4$
 & $k{=}4$: $2376/2400$ (three runs)\\
dodecahedron ($7776$)
 & $36$ & no cover below $\rho$ recorded
 & --- & star only, $\SCst\le35$
 & $k{=}3$: $2602/7776$\\
Cs\'asz\'ar torus ($1176$)
 & $11$ & no annealer cover recorded ($k{=}11$ cold: $563/1176$)
 & --- & star only, $\SCst\le10$ & ---\\
\bottomrule
\end{tabular}
\caption{First covers, descents and first stalls on the surfaces.  Cold
runs start from random sweeps, star-seeded runs from
Algorithm~\ref{alg:star}.  Sources:
\texttt{results/rp2\_k11\{,\_cont\}}\,\dots\,\texttt{rp2\_k6},
\texttt{hemicube\_k20}\,\dots\,\texttt{hemicube\_k7},
\texttt{klein\_k15\_8233\{,\_v2\}}\,\dots\,\texttt{klein\_k10},
\texttt{hemi\_octa\_k8}, \texttt{hemi\_dodeca\_k14}\,\dots\,\texttt{k7},
\texttt{icosa\_k23}\,\dots\,\texttt{icosa\_k5\_v2}; stalls in
\texttt{rp2\_k5\_*}, \texttt{hemicube\_k6\{,\_v2\}},
\texttt{klein\_k9\{,\_v2,\_v3,\_v4\}}, \texttt{hemi\_dodeca\_k6\{,\_v2\}},
\texttt{icosa\_k4\{,\_v2,\_v3\}}, \texttt{dodeca\_*}; rows marked
``(log)'' have only \texttt{logs/*.log} files.  Total wall is the sum of the
\texttt{elapsed\_s} of the rungs listed.  A stall is the best coverage in the
recorded runs of that $k$; it is not an impossibility proof.}
\label{tab:descent}
\end{table}

\subsection{The pipeline}

Putting the pieces together gives the recommended order of attack,
Algorithm~\ref{alg:pipeline}.  The third stage is \emph{only if
possible}: the CP-SAT model grows linearly in the number of facets of
$K\times K$ and in the chain length, and it exhausts memory at the scales
at which the annealer stalls (Section~\ref{sec:sat}).

\begin{algorithm}[H]
\caption{\textsc{Pipeline}($K$): the recommended order of attack on the
minimum cover of $K\times K$.}
\label{alg:pipeline}
\begin{algorithmic}[1]
\State \textbf{Star cover.}  Run Algorithm~\ref{alg:star}.  Record
$\rho(K)$, the certified bound $\SCst(K)\le\rho(K)-1$ and the seed
chains.  \textbf{If} a topological lower bound ($\TC(\|K\|)$) equals
$\rho(K)-1$, \Return: the bound is the minimum (even spheres).
\State \textbf{Annealer.}  $k\gets\rho(K)-1$; seed Algorithm~\ref{alg:anneal}
from the largest $k$ star domains.  \textbf{While} a complete verified
$k$-cover is found: record it, set $k\gets k-1$ and seed from the
incumbent.  Stop at the first $k$ where the coverage plateaus under the
guided ladder (Algorithm~\ref{alg:ladder}) for the available budget.
\State \textbf{CP-SAT.}  \textbf{If} the instance is small enough and
$k$ exceeds the topological floor, decide the $k$-domain instance for a
set of chain-length bounds $L$ (Section~\ref{sec:sat}).  A satisfying
assignment is a new certificate to be audited; an infeasibility result
bounds the invariant from below within the stated class of chains.
\end{algorithmic}
\end{algorithm}

A stall in step 2 is evidence and not a proof: the M\"obius band stalls at
$148/150$ for $k=3$ and the $3$-domain cover exists \cite{PM26}; the first
icosahedron run at $k=5$ stalled at $2395/2400$ for $200$ epochs and
completed under a warm continuation (Section~\ref{sec:icosa}).  Step 3 is
how a stall is converted into a statement, where the model allows.

\section{Case study: the icosahedron boundary}\label{sec:icosa}

The icosahedron boundary has $12$ vertices and $20$ triangles; $K\times K$
has $400\cdot6=2400$ facets.  It is a triangulation of $S^2$, so
$2\le\SCst\le\rho-1$ with $\TC(S^2)=2$ as the floor, and it is the
exhibit on which the three stages fit together best.

\emph{Stage 1.}  Lemma~\ref{lem:rows} gives $\rho\ge\lceil6\cdot20/5\rceil=24$
and the exact set cover attains it: $24$ rectangles, with domains of
$150\ (\times3)$, $126\ (\times5)$, $102\ (\times6)$, $78\ (\times7)$ and
$54\ (\times3)$ facets and chains of $2$ to $5$ links
(\texttt{results/icosa\_starcover\_opt}).  This certifies
$\SCst(\text{icosahedron})\le23$, in seconds.

\emph{Stage 2.}  Seeded from those chains (\texttt{--init}), the
annealer closed every rung from $k=23$ down to $k=6$, and then $k=5$
(Table~\ref{tab:icosa}).  The first $k=5$ run stalled at $2395/2400$
for all $200$ recorded epochs ($10\,270$\,s); a warm continuation of it
(\texttt{icosa\_k5\_v2}) completed the cover at its $67$th epoch
($2524$\,s), with domains of $894+755+387+185+179$ facets and chains of
$7$, $7$, $8$, $6$, $6$ links.  This is a verified $5$-domain
certificate, hence
\[
\SCst(\text{icosahedron})\le4,\qquad\text{so}\qquad 2\le\SCst(\text{icosahedron})\le4 .
\]
The annealer thus beat the product family it was seeded from by nineteen
domains, and closed to within two of the topological floor.  The wall
times of the twenty runs of Table~\ref{tab:icosa} sum to $17.9$ hours.
All twenty certificates, the $24$-domain seed and every rung $k=23,\dots,6$
as well as \texttt{icosa\_k5\_v2}, pass the independent auditor
(Section~\ref{sec:repro}).

\begin{table}[H]
\centering\footnotesize
\setlength{\tabcolsep}{4pt}
\begin{tabular}{@{}rrrr@{\hspace{18pt}}rrrr@{}}
\toprule
$k$ & epochs & min & coverage & $k$ & epochs & min & coverage \\
\midrule
$23$ & 23 & 26 & $2400$ & $13$ & 45 & 38 & $2400$ \\
$22$ & 67 & 81 & $2400$ & $12$ & 25 & 18 & $2400$ \\
$21$ & 99 & 127 & $2400$ & $11$ & 58 & 53 & $2400$ \\
$20$ & 2 & 3 & $2400$ & $10$ & 29 & 21 & $2400$ \\
$19$ & 69 & 69 & $2400$ & $9$ & 75 & 61 & $2400$ \\
$18$ & 93 & 118 & $2400$ & $8$ & 71 & 59 & $2400$ \\
$17$ & 2 & 2 & $2400$ & $7$ & 52 & 38 & $2400$ \\
$16$ & 31 & 29 & $2400$ & $6$ & 126 & 98 & $2400$ \\
$15$ & 22 & 18 & $2400$ & $5$ & 200 & 171 & $2395$ \\
$14$ & 3 & 3 & $2400$ & $5$ (cont.) & 67 & 42 & $2400$ \\
\bottomrule
\end{tabular}
\caption{The star-seeded descent on the icosahedron boundary ($2400$
facets): recorded epochs, wall time in minutes and final coverage of each
rung, from \texttt{results/icosa\_k23} through \texttt{icosa\_k5\_v2};
each rung is seeded from the previous one, the first from
\texttt{icosa\_starcover\_opt} ($24$ domains).  The row $k=5$ is the stall
at $2395/2400$; the row ``cont.'' is its continuation, the verified
$5$-domain record.  All runs use seed $8233$.}
\label{tab:icosa}
\end{table}

\emph{The stall at $k=4$.}  Three runs target four domains
(\texttt{icosa\_k4}, \texttt{icosa\_k4\_v2}, \texttt{icosa\_k4\_v3}), the
first seeded from the four largest domains of the $5$-domain record and
each later one from its predecessor.  Each ran $200$ recorded epochs
($9572$, $8590$ and $8635$\,s) and each ends at exactly $2376/2400$
facets, $24$ short of a cover.  The first reached $2376$ at epoch $160$;
the other two start there and do not improve.  They share a seed and a
lineage, so they are one observation repeated, not three independent
ones, and the lesson of the $k=5$ stall---which a continuation did
complete---is that a plateau of this kind is evidence and not proof.  We
therefore leave $2\le\SCst\le4$ as it stands: whether a four-domain cover
exists is open.  No CP-SAT run on this complex is recorded; at $2400$
facets, a model of the type described in Section~\ref{sec:sat} is far
beyond the $500$-facet $\partial\Delta^4$ model, which already needs half
an hour at $L=48$.

This is the concrete demonstration that star cover followed by annealer
is much easier than the cold, high-$k$ start of Section~\ref{sec:descent}:
the pipeline reached $\SCst\le4$ on a $2400$-facet complex without ever
running a cold search.

\section{Exact decisions by CP-SAT}\label{sec:sat}

When the annealer stalls at a $k$ above the topological floor, we can
ask a solver.  The decision problem is: does $K\times K$ admit a cover
by at most $k$ contiguity domains whose chains have at most $L$ links?
Our encoding (\texttt{paper-overleaf/paper/scs\_sat.py}) is for the
CP-SAT solver of OR-Tools \cite{ORTools}.  For each of the $k$ domains
it has one vertex-valued variable per grid vertex and per interior map
of a chain $\pi_1=\varphi_0,\dots,\varphi_L=\pi_2$; a facet is covered
by a domain when, for each link, the union of the images of the facet
under the two maps lies in a maximal facet of $K$; every facet must be
covered by some domain.  A chain with fewer than $L$ links is padded by
repeating $\pi_2$, and padding never shrinks a domain, so ``no cover with
chains of exactly $L$ links'' is equivalent to ``no cover with chains of
at most $L$ links'' and infeasibility is monotone in $L$.  For
$\partial\Delta^4$ and $k=2$ the model has
$5000(L-2)+50(L-1)+1005$ variables.

\subsection{Validating the instrument}

An infeasibility answer from a hand-built model is worth only as much as
the model.  The encoding was audited before any result was used
(\texttt{experiments/}\allowbreak\texttt{s3\_sat\_validation\_log.md}; scripts in
\texttt{scripts/}\allowbreak\texttt{sat\_validation/}).
(A1)~On $1100$ random chains ($500$ on $\partial\Delta^4$, $300$ on the
circle, $300$ on the Klein bottle; $L\in\{2,\dots,7\}$) the encoding's
notion of ``chain covers facet'' agrees with the definition of
Section~\ref{sec:back} and with the annealer's own
\texttt{chain\_bad\_counts} on all three: $1100/1100$.
(A2)~All five recorded $3$-domain covers of $\partial\Delta^4$ round-trip
through the encoding (every recorded domain is contained in the
recomputed one, and padding to $16$ links never shrinks a domain).
(A3)~Solver controls in both directions: $\partial\Delta^4$ with
$k=3$, $L=4$ is satisfiable; the circle with $k=1$ is infeasible and with
$k=2$ satisfiable; the M\"obius band with $k=2$, $L=6$ is infeasible and
with $k=3$ satisfiable---five of five.  (A4)~The audit found and fixed
three defects of the tool (no guard for $L<2$, misleading labels, and a
truncation of warm chains in hybrid mode), none of which touches the
cold-start verdicts reported here.  On 3~October~2026 we re-ran the scripts of (A1) and (A2) and the
five controls of (A3), and the rungs $L=4,5,8,16$ of the ladder below, on a
fresh installation of OR-Tools~9.15; all outcomes reproduced.

\subsection{The ladder on \texorpdfstring{$\partial\Delta^4$}{the 3-sphere model}}

\begin{table}[H]
\centering\footnotesize
\begin{tabular}{@{}rrrlr@{}}
\toprule
$L$ & variables & constraints & result for $k=2$ & recorded time\,(s) \\
\midrule
$4$ & $11\,155$ & $156\,501$ & infeasible & $12$\\
$5$ & $16\,205$ & $227\,501$ & infeasible & $33$\\
$8$ & $31\,355$ & $440\,501$ & infeasible & $77$\\
$16$ & $71\,755$ & $1\,008\,501$ & infeasible & $239$\\
$24$ & $112\,155$ & $1\,576\,501$ & infeasible & $799$\\
$32$ & $152\,555$ & $2\,144\,501$ & infeasible & $983$\\
$40$ & $192\,955$ & $2\,712\,501$ & infeasible & $1622$\\
$48$ & $233\,355$ & $3\,280\,501$ & infeasible & $1988$\\
\bottomrule
\end{tabular}
\caption{CP-SAT on $\partial\Delta^4\times\partial\Delta^4$ ($500$
facets), $k=2$.  The logs of $L=24,\dots,48$ are in
\texttt{logs/sat\_ladder/}; those of $L=4,5,8,16$ are recorded in
\texttt{experiments/s3\_sat\_validation\_log.md}; a further smoke run
($L=12$, infeasible in $60$\,s) is \texttt{results/s3\_k2\_sat\_smoke\_L12}.
The solver reports no proof log.}
\label{tab:sat}
\end{table}

\begin{proposition}[solver result, bounded chains]\label{prop:sat}
There is no cover of $\partial\Delta^4\times\partial\Delta^4$ by two
contiguity domains whose chains have at most $48$ links.
\end{proposition}

We label Proposition~\ref{prop:sat} as a \emph{solver result}: the
encoding is validated as above, but the infeasibility is the answer of a
solver that returns no checkable certificate (its log reads
\texttt{LRAT\_status: NA}), and we have not re-checked it by any
independent method.  It is also a statement about a \emph{bounded}
class of covers, and that is the point on which this section is
careful.  The parameter $L$ counts links of a chain, not chains or
domains: the number of domains is $k=2$ throughout.

\subsection{What would make it unconditional}

To pass from Proposition~\ref{prop:sat} to $\SCst(\partial\Delta^4)\ne1$
one needs the compression claim~C2 of Conjecture~\ref{conj:compress}:
every contiguity domain admits a chain of at most $48$ links.  What we
have is evidence, not a proof.  On $\partial\Delta^4$ two maps are
$1$-contiguous on a product facet exactly when their images on its seven
grid vertices use at most four distinct values, from which one proves
(a)~an \emph{interpolation lemma}: if $\varphi\sim\psi$ on $D$, every
map on the Hamming box between them is contiguous to both; and
(b)~every single product facet has a two-link chain through a constant.
The naive route from these to a bound---walk from $\pi_1$ to $\pi_2$
one grid vertex at a time---is \emph{falsified}: $81\,820$ pairs of
facets force cycles in the order of the flips, so a short chain, where
one exists, uses intermediate maps outside the interpolation box
(constants, value-dropping maps).  Measured, using the same constraints
at $k=1$ to compute the exact minimum number of links of a given domain:
the star-cover domains need $2$ links, the three-domain annealer covers
$3$ links, the staircase rectangles $2$ or $3$, and the domains of
six $k=2$ plateau runs---the largest chainable sets the annealer ever
found---$2$ to $4$; a $320$-facet domain rejected all $131$ attempted
additions; random sets of $400$ or more facets are not chainable at all
within $24$ links.  The maximum is $4$, twelve times below $48$.  We
state this as a conjecture and do not attempt a proof here.

\subsection{Consequences}

If C2 holds, then Proposition~\ref{prop:sat} gives
$\SCst(\partial\Delta^4)=2$ at level $0$ while $\TC(S^3)=1$.  That would
be a \emph{second} example of a genuine level-$0$ gap between the
topological invariant and its discretization, which vanishes only in the
barycentric limit \cite{G18}.  The first is the M\"obius band,
$\SCst(K_M)=2>1=\TC$ \cite{PM26}, itself conditional on chains of at
most $8$ links; there the annealer stalls at $148/150$ and CP-SAT
produces the missing $3$-domain cover (\texttt{results/mobius\_k3\_sat},
chains of $5$--$6$ links, which passes the independent auditor) and
proves $k=2$ infeasible (we reproduced $L=6$ in $10.5$\,s).  A gap on
$\partial\Delta^4$ would move the phenomenon from a band with boundary to
a sphere.  All of this is conditional on C2.

\subsection{Limits of the method}

The model does not scale to the instances where the annealer stalls.
For $\partial\Delta^6$ ($S^5$) and $k=2$ the models for $L=3$ and $L=4$
($197\,771$ and $370\,741$ variables) were built, using about $3$\,GB,
and the solver then failed with \texttt{std::bad\_alloc} under a $5$\,GB
cap on a $14$\,GB machine.  For $\mathbb{RP}^2$ with $k=5$ and $L=4$
($63\,546$ variables) a run was killed by memory exhaustion at the $6$\,GB
cap after $4$\,h $35$\,min without reporting a solution or a bound.  No
verdict was obtained for $\mathbb{RP}^2$ at $k=5$ or for the Klein bottle at
$k=9$.  CP-SAT is therefore a stage for small complexes: the M\"obius
band ($150$ facets) and the $3$-sphere model ($500$ facets).  This is
why the pipeline of Algorithm~\ref{alg:pipeline} uses it ``only if
possible''.

\section{Results}\label{sec:results}

Table~\ref{tab:records} places the new results alongside the earlier
bounds.  The ``before'' column accumulates the
published pipeline \cite{OLGB21} and its reimplementations in the first
phase of this project (Section~\ref{sec:intro}): the $96$-core C++/MPI
version and the faithful single-GPU port.  The ``after'' column is
the chain-first annealer of Section~\ref{sec:annealer}.  The spheres,
the circle, the wedge, the cube boundary and the M\"obius band are in
Table~\ref{tab:records}; the surface campaigns that rest on the pipeline
of Algorithm~\ref{alg:pipeline} are in Table~\ref{tab:surfaces}.

\begin{remark}
The annealer certifies upper bounds, not optima.  A successful $k$-chain
run proves $\SCst(K)\le k-1$ but does not prove that $k-1$ is the
minimum; we improve on the bounds obtained in the recorded runs of the
algorithms of \cite{OLGB21}, not claiming that every reported value is a global
minimum.  The sharp equalities in the tables arise from independent
lower bounds---the values $\TC(S^n)=2$ for even $n\ge2$ and
$\TC(S^1)=1$ \cite{F03}, and $\TC(S^1\vee S^1)\ge2$ (cohomological lower bound,
\cite{F03}) for the wedge---matched by the upper bounds of our certificates (for the
circle and the wedge, also by the earlier constructions of
\cite{OLGB21}).  A construction gives an upper bound; equality needs
the lower bound.  For spaces not yet
pinned by a matching lower bound the reported values are the best known
upper bounds.
\end{remark}

\begin{table}[H]
\centering\footnotesize
\begin{tabular}{@{}>{\raggedright\arraybackslash}p{2.7cm}>{\raggedright\arraybackslash}p{3.5cm}>{\raggedright\arraybackslash}p{2.8cm}>{\raggedright\arraybackslash}p{4.2cm}@{}}
\toprule
target & before & after & certificate \\
\midrule
circle $\partial\Delta^2$ ($18$)
  & $57$\,s \cite{OLGB21}; level $1\approx10$\,h
  & $<15$\,s
  & $9+9$, $\SCst\le1$, so $\SCst=1=\TC$, audited\\
wedge $K_4\setminus e$ ($50$)
  & $25$\,min \cite{OLGB21}, $3$ domains
  & $41$\,s
  & $18+16+16$, $\SCst=2$, audited\\
cube boundary ($864$)
  & $\SCst\le 7$ (product bound, $\rho=8$)
  & $\approx 24$\,min total; $304$\,s final guided stage
  & $482+236+146$, $\SCst=2=\TC$, audited\\
M\"obius band $K_M$ ($150$)
  & not attempted
  & $17$\,s (warm from $k{=}5$)
  & $54+37+34+25$, $\SCst\le3$, audited$^{\ddagger\ddagger}$\\
$\partial\Delta^3$ (S$^2$, $96$)
  & never on $96$ cores, $120$\,h jobs;
    $4.4$\,min on T4 (first phase)
  & $22$\,s
  & $54+29+13$, $\SCst=2=\TC$, audited\\
$\partial\Delta^4$ (S$^3$, $500$)
  & $\SCst\le 19$ (initial GPU baseline, $20$ domains, $\approx3$\,h GPU, first phase)$^{*}$
  & $153$\,s, one epoch
  & $309+130+61$, $\SCst\le 2$, audited\\
$\partial\Delta^5$ (S$^4$, $2520$)
  & not attempted
  & $13$\,min, one epoch
  & $1634+653+233$, $\SCst=2=\TC$, audited\\
$\partial\Delta^6$ (S$^5$, $12348$)
  & not attempted
  & $34$\,min, epoch $2$
  & $7061+4239+1048$, $\SCst\le 2$, audited\\
$\partial\Delta^7$ (S$^6$, $59136$)
  & not attempted
  & ${\approx}2.5$\,h, $8$-rung ladder$^\S$
  & $44909+11451+2776$, $\SCst=2=\TC$, audited\\
all $\partial\Delta^{n+1}$
  & S$^1$, S$^2$, S$^3$ case by case
  & $\SCst\le2$ for all $n\ge1$; $=\TC$ for even $n\ge2$
  & Corollary~\ref{cor:spheres}\\
\bottomrule
\end{tabular}
\caption{Then and now.  Facet counts of the verified partition are listed
in the certificate column.  Every certificate in the last column was
re-audited independently of the search code.  $^{*}$A later run under a batched-retry configuration
(\texttt{results/s3\_20260920\_072055/seed17},
commit \texttt{6857ce4} of the GPU repository, about $350$~minutes)
produced an independently checked $13$-domain cover,
$\SCst(\partial\Delta^4)\le12$; listed for completeness, it does not
improve the three-domain certificate.
$^{\ddagger\ddagger}$A companion note \cite{PM26} reports
$\SCst(K_M)=2$ (conditional on chains $\le8$ links) via
constraint programming (OR-Tools CP-SAT), exhibiting a subdivision gap
of one unit: $\TC(\|K_M\|)=\TC(S^1)=1$ while $\SCst(K_M)=2$ at
level~$0$.}
\label{tab:records}
\end{table}

\begin{table}[H]
\centering\footnotesize
\setlength{\tabcolsep}{3pt}
\begin{tabular}{@{}lrrrr>{\raggedright\arraybackslash}p{4.1cm}c>{\raggedright\arraybackslash}p{1.9cm}@{}}
\toprule
complex & $|V|$ & $|K{\times}K|$ & $\rho$ & $\rho{-}1$ &
best certified $\SCst\le$ (source) & $\TC$ & independent audit \\
\midrule
M\"obius band & $5$ & $150$ & $4$ & $3$
 & $2$ (CP-SAT, $3$ domains, \texttt{mobius\_k3\_sat}; $=2$ within chains $\le8$ \cite{PM26})
 & $1$ & yes\\
hemi-icosahedron & $6$ & $600$ & $6$ & $5$
 & $5$ (annealer $k{=}6$, \texttt{rp2\_k6}); $k{=}5$ open
 & $3$\,\cite{FTY03} & yes\\
hemicube & $7$ & $864$ & $7$ & $6$
 & $6$ (\texttt{hemicube\_k7}); $k{=}6$ stalls at $836/864$
 & $3$ & not supported\\
hemi-octahedron & $9$ & $1536$ & $9$ & $8$
 & $7$ (\texttt{hemi\_octa\_k8}, star-seeded)
 & $3$ & not supported\\
hemi-dodecahedron & $10$ & $1944$ & $15$ & $14$
 & $6$ (\texttt{hemi\_dodeca\_k7}, star-seeded); $k{=}6$ stalls at $1922/1944$
 & $3$ & not supported\\
Klein bottle & $8$ & $1536$ & $12$ & $11$
 & $9$ (\texttt{klein\_k10}); $k{=}9$ stalls at $1494/1536$
 & $4$\,\cite{CV17} & yes\\
Cs\'asz\'ar torus & $7$ & $1176$ & $11$ & $10$
 & $10$ (star cover only)
 & $2$\,\cite{F03} & yes (star)\\
cube boundary & $8$ & $864$ & $8$ & $7$
 & $2$ (annealer, $k{=}3$); $=\TC$
 & $2$ & yes\\
icosahedron boundary & $12$ & $2400$ & $24$ & $23$
 & $4$ (star-seeded descent, \texttt{icosa\_k5\_v2}); $k{=}4$ stalls at $2376/2400$
 & $2$ & yes (ladder)\\
dodecahedron boundary & $20$ & $7776$ & $36$ & $35$
 & $35$ (star cover only); $k{=}3$ stalls at $2602/7776$
 & $2$ & yes (star)\\
\bottomrule
\end{tabular}
\caption{The surface campaigns: star-rectangle constant $\rho$, the
star bound $\rho-1$, and the best certified bound on $\SCst$ in the
repository, with the topological lower bound $\TC$ (reduced convention;
$\SCst\ge\TC$).  The M\"obius band, the cube and the icosahedron are
the rows on which the annealer or the solver certifies less than the star
bound by a wide margin.  ``Independent audit'' refers to the audit
re-run for this paper (Section~\ref{sec:repro}), not to the
\texttt{certification} field stored in the older \texttt{summary.json}
files; the hemicube, hemi-octahedron and hemi-dodecahedron certificates
pass only the internal verifier.  The row for the Klein bottle reports
the audited star cover (\texttt{klein\_starcover\_greedy}, $12$
domains, here $\rho=12$) together with the audited rungs
\texttt{klein\_k10}\,\dots\,\texttt{klein\_k14} and \texttt{klein\_k15\_8233\_v2}.}
\label{tab:surfaces}
\end{table}

\emph{The surface campaigns, read with the star bound.}  Table~\ref{tab:surfaces}
shows the three behaviours the pipeline should be expected to meet.  On
the hemi-icosahedron the annealer's best certificate equals the star
bound: the descent from $k=11$ went exactly to $\rho$ and stalled one
below, so the star cover would have given the same certificate in
seconds.  On the hemicube, the hemi-octahedron and the
hemi-dodecahedron, the annealer beat the star bound by $0$, $1$ and $8$
respectively (certified $6$, $7$ and $6$ against $\rho-1=6$, $8$ and
$14$), and the dodecahedron boundary and the torus keep the star bound as
their only certificate.  On the cube, the Klein bottle and the
icosahedron the margins are $5$, $2$ and $19$.  Lower bounds
are far from these: the invariant of every row is between $\TC$ and the
certified value, and only the cube (and the spheres of
Table~\ref{tab:annealer}) is pinned.

Several entries deserve emphasis.

\emph{The wedge and the cube separate free-form search from product
covers.}  On the wedge of two circles the product bound gives only
$\SCst\le 3$ ($\rho=4$), while the annealer reproduces the sharp
certificate of \cite{OLGB21} in $41$\,seconds---finding a different
partition ($18+16+16$ vs.\ $20+18+12$) with shorter chains ($2$, $5$,
$5$ links vs.\ $8$, $11$, $10$).  On the cube boundary the gap is
more dramatic: $\rho=8$ forces any cover by the star rectangles $D(v,w)$ defined in
Theorem~\ref{thm:prod} to use at least $8$ domains, certifying only $\SCst\le7$, yet the chain-first annealer
covers all $864$ facets with just $3$ domains and confirms
$\SCst=2=\TC(S^2)$.  In neither case is the optimal number of domains attained by the
star-rectangle construction of Theorem~\ref{thm:prod}.

\emph{The $4$-sphere and $5$-sphere were priced out of the old world.}
$\partial\Delta^5\times\partial\Delta^5$ has $2520$ maximal facets and
$\partial\Delta^6\times\partial\Delta^6$ has $12348$; neither was
attempted with the greedy pipeline.  The chain-first annealer solves
them in $13$ and $34$~minutes, respectively, and delivers independently
audited certificates.  Level-$0$ sharpness
$\SCst(\partial\Delta^5)=\TC(S^4)=2$ follows from the certificate
together with Farber's lower bound $\TC(S^4)=2$.

\emph{The $6$-sphere pushes the product further.}
$\partial\Delta^7\times\partial\Delta^7$ has $59\,136$ maximal
facets---nearly five times the $5$-sphere model.  An eight-rung ladder
campaign on an NVIDIA RTX PRO 6000 Blackwell Server Edition covers all
$59\,136$ facets in approximately $2.5$ hours: a cold five-rung phase
(population $S{=}64$) reaches $99.89\%$ coverage ($59\,072/59\,136$),
and a guided three-rung continuation at $4\times$ population ($S{=}256$)
closes the remaining $64$-facet gap.  The certificate gives
$\SCst(\partial\Delta^7)=\TC(S^6)=2$, another sharp even-dimensional
value.  With the addition of $S^6$, the annealer has produced explicit
certificates for every boundary-of-simplex sphere from $S^1$ through
$S^6$---a facet-count range spanning three orders of magnitude ($18$ to
$59\,136$).

\emph{Prediction became quantitative.}  Before the corrected runs,
rectangle arithmetic predicted two-chain plateaus at $460=23\cdot 20$
facets out of $500$ on $\partial\Delta^4$ and at $2380=34\cdot 70$ out
of $2520$ on $\partial\Delta^5$.  Both coverage values were observed
exactly, the latter in a single epoch.  The agreement is at the level of
the numbers: the domains returned at those values are often not
rectangles (Section~\ref{sec:annealer}), so what the plateau measures is
how far local dynamics reach, not that a two-chain optimum has product
structure.

\emph{The search and the theorem certify the same optimum.}  The
$2$-sphere and $4$-sphere values match $\TC(S^2)=\TC(S^4)=2$ given by
Farber's theorem \cite{F03}, independently certifying sharpness at
subdivision level~$0$.

\emph{The M\"obius band shows both the annealer's reach and its limits
on small complexes.}  The vertex-minimal triangulation $K_M$ of the
M\"obius band ($5$ vertices, $5$ triangles, $150$ product facets) has
$\TC(\|K_M\|)=\TC(S^1)=1$, so $\SCst(K_M)\ge1$.  A cold $k{=}5$
search covers all $150$ facets in $36$~seconds; warm-starting from the
certified $5$-domain cover, a $k{=}4$ search covers them in $17$~seconds
on the first epoch, certifying $\SCst(K_M)\le3$.  Eight independent
$k{=}3$ attempts---cold, warm-restarted from the $4$-domain certificate,
and guided---all stalled at exactly $148$ of $150$ facets.  A companion
note \cite{PM26} closes the gap to $\SCst(K_M)=2$ (conditional on chain
length $\le8$) using a complete constraint-programming search (OR-Tools
CP-SAT), which also proves that no $2$-domain system with chains of up
to $8$ links exists.  The M\"obius band at $150$ facets and the
$3$-sphere model at $500$ facets are in reach of exact methods
(Section~\ref{sec:sat}); the larger models are not.

\section{Certificates, auditing, reproduction}\label{sec:repro}

Randomized search output is usually taken on trust; none of the numbers
above is.  Every reported system is discharged as data: the winning
state's chains are replayed deterministically (identical splitmix64 RNG
on GPU and CPU), each chain is shortened by the greedy reduction of
\cite[Algorithm 2]{OLGB21}, and the system is asserted to connect the
two projections through pairwise $1$-contiguous simplicial maps over
every facet of every domain.  An independent auditor---a pure-Python
program that rebuilds each triangulation from scratch and shares no
code with the search---re-runs every check; it has caught real bugs
during development and exits nonzero on any discrepancy.

We distinguish four operations, which should not be conflated.
(i)~\emph{Verifying saved files}: the pinned auditor re-checks archived
covers and planner maps against a hash manifest; on 22~September~2026
this verified eight certificates ($S^2$, $S^3$, $S^4$, $S^5$, $S^6$, the cube
boundary, the wedge and the M\"obius band).  (ii)~\emph{Reproducing an explicit CPU
construction}: a reproduction script builds explicit product
certificates for the circle, $S^2$, $S^3$ and $S^4$ and re-audits them;
the figure of about two minutes on CPU has not been re-measured for
this whole scope.  (iii)~\emph{Repeating a GPU search}: an optional GPU
mode replays the flagship annealer runs.  (iv)~\emph{Reproducing
timings}: the wall times reported here ($22$\,s, $153$\,s, $13$\,min,
\dots) are recorded single-run observations.  A certificate audit
(i) does not by itself validate (ii)--(iv).  Certificate documents listing every domain's
facets and every planner map are generated from data by a single
generator program.  The tables of the first phase are typed from audited
output; the new tables of Sections~\ref{sec:products}--\ref{sec:sat} are
read from the \texttt{summary.json} files and logs named in their captions,
and the values of Table~\ref{tab:rho} were recomputed for this paper.  A
fifth category is not a certificate at all: the infeasibility results of
Section~\ref{sec:sat} are solver outputs, validated at the level of the
encoding but not independently re-checked.

\paragraph{Audit status of the star covers and ladders.}  On
3~October~2026 we ran the independent auditor
(\texttt{src/independent\_check.py}) on the certificates of
Sections~\ref{sec:products}--\ref{sec:icosa}.  It verified (exit status
zero) the star covers of the circle, of $\partial\Delta^{n+1}$ for
$n=2,\dots,8$ (the last with $1\,287\,000$ facets), of the cube boundary,
the M\"obius band, the hemi-icosahedron, the Klein bottle, the
Cs\'asz\'ar torus (greedy and exact), the icosahedron and the
dodecahedron; the entire icosahedron ladder (the $24$-domain seed, the
rungs $k=23,\dots,6$ and the $5$-domain record \texttt{icosa\_k5\_v2}:
twenty certificates); the hemi-icosahedron rungs $k=6,7$; the Klein
rungs $k=10,\dots,15$ that exist as certificates; the M\"obius
certificates \texttt{mobius\_k3\_sat} and \texttt{mobius\_k4\_from\_k5};
and the earlier records \texttt{s3\_k3\_v2}, \texttt{cube\_k3\_guided}
and \texttt{circle\_probe3}.  It does not support the hemicube, the
hemi-octahedron or the hemi-dodecahedron: it selects the complex from the
facet count, the counts $864$ and $1536$ coincide with those of the cube
and the Klein bottle, and $1944$ is unsupported.  The certificates
\texttt{hemicube\_k7}, \texttt{hemi\_octa\_k8} and
\texttt{hemi\_dodeca\_k7}, their star covers and the remaining rungs of
those three descents therefore pass only the internal verifier, which
re-checks every map and link with the search's own data structures; we
report them as such (Table~\ref{tab:surfaces}).  The
\texttt{independent\_audit} field of the older \texttt{summary.json}
files still reads \texttt{not\_run}, because the audit was run outside
the export step.

\paragraph{Repositories.}  The project is split across two public repositories.
The \emph{source code}---the chain-first annealer, the certificate builder,
the independent auditor, and the reproduction script---lives at
\begin{center}
\url{https://github.com/MiguelPerezI/chain-first-annealer}
\end{center}
The \emph{certificates} of eight certified systems ($S^2$, $S^3$, $S^4$,
$S^5$, $S^6$, the cube boundary, the wedge and the M\"obius band)---every domain and every planner
map (as \texttt{cover.txt}, \texttt{planners.txt}
and \texttt{summary.json} per run, and as readable certificate documents)---are
archived separately at
\begin{center}
\url{https://github.com/MiguelPerezI/simplicial-complexity-certificates}
\end{center}
The circle certificate
(\texttt{chain-first-annealer/opti/results/circle\_probe3}) and the
partial plateau states used in Section~\ref{sec:annealer} are not in the
certificates archive; they are in \texttt{chain-first-annealer/opti/results/}.
A reader who wishes to verify, say, the $4$-sphere cover by hand can
open \texttt{cover.txt} and \texttt{planners.txt} in
\texttt{results/chain-first-annealer/s4\_k3\_v2/} of the certificates
repository, or run the auditor \texttt{opti/src/independent\_check.py}
of the source repository, from the directory that contains both clones:
{\footnotesize
\begin{verbatim}
python3 chain-first-annealer/opti/src/independent_check.py \
  simplicial-complexity-certificates/results/chain-first-annealer/s4_k3_v2
\end{verbatim}}
\noindent The whole archive can be checked with the pinned auditor and
hash manifest, from the root of \texttt{simplicial-complexity-certificates}:
{\footnotesize
\begin{verbatim}
python3 audit/verify_all.py --check
\end{verbatim}}

\paragraph{Pinned commits.}  The results of the first phase (Sections~\ref{sec:annealer}
and~\ref{sec:results}, Table~\ref{tab:records}) were produced at
commit \texttt{9e262d7} of the source repository and commit
\texttt{e2bdb5b} of the certificates repository.  The star covers, the
descents, the icosahedron case study and the CP-SAT material of
Sections~\ref{sec:products}--\ref{sec:sat} postdate those commits: they
are in \texttt{opti/results/}, \texttt{opti/logs/} and
\texttt{opti/experiments/} of the source repository at commit
\texttt{7271048}, and the new certificates are not yet in the certificates
archive.
% TODO(authors): archive the star-cover, ladder and mobius_k3_sat
% certificates in simplicial-complexity-certificates and re-pin both commits.

\paragraph{Hardware and software.}  All annealer runs except $S^6$ were
executed on a single NVIDIA Tesla T4 GPU\@.  The $S^6$ campaign ran on a
single NVIDIA RTX PRO 6000 Blackwell Server Edition ($24$\,GB, compute
capability $12.0$).  Wall times are single-run observations at a single
seed per target (Table~\ref{tab:params}); they are not means over
repeated trials.  The surface campaigns of Section~\ref{sec:descent}
and the CP-SAT runs of Section~\ref{sec:sat} ran on other machines whose
hardware is not attested in the stored summaries; the exact-$\rho$
computations and the checks of this section ran on a four-thread laptop CPU.

\paragraph{Known issues.}  The reproduction script
\texttt{reproduce\_all.sh} uses \texttt{set -e} without
\texttt{pipefail}; piping auditor output into \texttt{tail} can mask a
nonzero exit status.  The exporter's \texttt{verify\_system} relies on
\texttt{assert}, which is stripped by \texttt{python -O}.  Neither
issue invalidates the published certificates, which were checked by
the independent auditor, but both should be addressed before the
scripts are treated as a standalone reproduction gate.

\section{Concluding remarks}\label{sec:concl}

The arc of this project is that the expensive object turned out to be
the wrong object, twice.  The randomized greedy hunts chains one facet at
a time, at a cost that defeated a $96$-core cluster; the annealer hunts
whole systems of chains and solves the classical targets in seconds to
hours; and the coverage values its searches kept returning pointed to a
theorem that makes the first move search-free.  What we now recommend is
the pipeline of Algorithm~\ref{alg:pipeline}: the star cover first,
because it is instant, unconditional, exact in its own family and a seed;
the annealer second, descending in $k$ below $\rho$; CP-SAT third, and
only where the model is small enough to decide.  The star cover is also
a floor under our own failures: for the dodecahedron boundary and the
Cs\'asz\'ar torus it is the only certified bound we have
(Table~\ref{tab:surfaces}), and for the hemi-dodecahedron it
supplied the seed from which every certificate in the sequence
$k=14,\dots,7$ descends, where cold starts at $k=15,20,30$ produced none.

Three kinds of statement appear in this paper and we have tried never to
conflate them.  \emph{Audited certificates} are explicit covers verified
link by link and, with the exceptions listed in Section~\ref{sec:repro},
re-checked by an independent program.  \emph{Solver results} are the
CP-SAT infeasibility statements of Section~\ref{sec:sat}; they hold for
a stated bound on the chain length, depend on a validated but
unproved-in-the-paper encoding, and are outputs of a solver whose
proofs we do not independently re-check.  \emph{Empirical observations}
are the plateaus and stalls: recorded runs, single seeds, finite
budgets.  A stall, such as $2376/2400$ for $k=4$ on the icosahedron,
$460/500$ for $k=2$ on $\partial\Delta^4$, $148/150$ for $k=3$ on the
M\"obius band, is evidence that a cover may not exist; the M\"obius
band is the standing reminder that it can exist after all \cite{PM26}.

The open questions are of three sizes.  On the icosahedron boundary,
$2\le\SCst\le4$: does a four-domain cover exist?  On
$\partial\Delta^4$, the two conjectures below.  And, for every row of
Table~\ref{tab:surfaces} below its star bound, whether the stalled value
is the minimum.

\begin{conjecture}\label{conj}
\emph{(C1)} $\SCst(\partial\Delta^4)=2$ at level $0$; equivalently,
with the convention $\SCst=(\text{minimum number of domains})-1$ used
throughout, two domains do not suffice on the S$^3$ model, and the
level-$0$ invariant of this model exceeds $\TC(S^3)=1$.
\end{conjecture}

\noindent\emph{Status of C1.}
\emph{Solver result, conditional on chain length (Section~\ref{sec:sat}).}
There is no cover of $\partial\Delta^4\times\partial\Delta^4$ by two
contiguity domains whose chains have at most $48$ links
(CP-SAT: infeasible at $L=4,5,8,16,24,32,40,48$; padding makes the
statement monotone in $L$).  C1 is therefore settled within that class
of covers, and the class contains every cover of $\partial\Delta^4$ we have
recorded (longest reduced chain: $9$ links).  \emph{Unconditionally C1 is open}:
it follows from the solver result together with the compression
conjecture below, and not from the solver result alone.
\emph{Empirical.}  Eight stored two-chain annealer searches on
$\partial\Delta^4$ (seeds $21$--$24$ at $\beta{=}0.05$; seed $5$ at
$\beta{=}0.05$, $\gamma{=}0$ and at $\beta{=}5$, $\gamma{=}0.5$; seeds
$53$ and $54$ at $\beta{=}5$, $\gamma{=}0.5$;\footnote{The eight runs are
\texttt{s3\_k2\_A} through \texttt{s3\_k2\_D},
\texttt{s3\_k2\_longcor\_5}, \texttt{s3\_k2\_longflat\_5},
\texttt{s3\_k2\_v2\_point} and \texttt{s3\_k2\_v2\_track}, all in
\texttt{chain-first-annealer/opti/results/}.} both old and corrected
energy) all plateau at coverage $460/500$ facets, and the largest
single two-chain domain in the stored runs has $320$ facets.  A failed
search is not an impossibility proof; this paragraph is the experimental
support that preceded the solver result, not part of it.

\begin{conjecture}[compression]\label{conj:compress}
\emph{(C2)} Every contiguity domain $D\subseteq\partial\Delta^4\times
\partial\Delta^4$ admits a contiguity chain from $\pi_1|_D$ to
$\pi_2|_D$ with at most $48$ links.
\end{conjecture}

\noindent C1 follows from C2 and the solver result of
Section~\ref{sec:sat}: a two-domain cover would have chains of at most
$48$ links.  Support for C2 is measured, not proved
(Section~\ref{sec:sat}): every chainable domain we measured needs at
most $4$ links, a twelvefold margin.  We state C2 as a conjecture and
do not attempt a proof here.

The plateau states listed for $\partial\Delta^4$ and
$\partial\Delta^5$ attain coverage values that coincide with those of
explicit two-star-rectangle constructions. This is an empirical statement
about the recorded runs, not a universal convergence claim or a geometric
classification of their domains. Some saved domains are contained in no
star rectangle. The circle admits a complete two-domain cover despite
$\rho(\partial\Delta^2)=3$.  On $\partial\Delta^3$, the sixteen greedy
replicas of Section~\ref{sec:fail} froze at near-rectangle fronts of
$30$--$54$ facets within their recorded budgets; these are
finite-budget observations, not further evidence for a general
convergence statement.

\begin{thebibliography}{10}

\bibitem{Cerny85}
V.~\v{C}ern\'y.
\newblock Thermodynamical approach to the traveling salesman problem:
  an efficient simulation algorithm.
\newblock {\em J. Optim. Theory Appl.}, 45(1):41--51, 1985.

\bibitem{CV17}
D.~C. Cohen and L.~Vandembroucq.
\newblock Topological complexity of the Klein bottle.
\newblock {\em J. Appl. Comput. Topol.}, 1(2):199--213, 2017.

\bibitem{F03}
M.~Farber.
\newblock Topological complexity of motion planning.
\newblock {\em Discrete Comput. Geom.}, 29(2):211--221, 2003.

\bibitem{FTY03}
M.~Farber, S.~Tabachnikov, and S.~Yuzvinsky.
\newblock Topological robotics: motion planning in projective spaces.
\newblock {\em Int. Math. Res. Not.}, 2003(34):1853--1870, 2003.

\bibitem{G18}
J.~Gonz\'alez.
\newblock Simplicial complexity: piecewise linear motion planning in
  robotics.
\newblock {\em New York J. Math.}, 24:279--292, 2018.
\newblock arXiv:1701.07612.

\bibitem{KGV83}
S.~Kirkpatrick, C.~D. Gelatt, Jr., and M.~P. Vecchi.
\newblock Optimization by simulated annealing.
\newblock {\em Science}, 220(4598):671--680, 1983.

\bibitem{Metropolis53}
N.~Metropolis, A.~W. Rosenbluth, M.~N. Rosenbluth, A.~H. Teller, and
  E.~Teller.
\newblock Equation of state calculations by fast computing machines.
\newblock {\em J. Chem. Phys.}, 21(6):1087--1092, 1953.

\bibitem{OLGB21}
C.~Ortiz, A.~Lara, J.~Gonz\'alez, and A.~Borat.
\newblock A randomized greedy algorithm for piecewise linear motion
  planning.
\newblock {\em Mathematics}, 9(19):2358, 2021.
\newblock arXiv:2008.13290.

\bibitem{ORTools}
L.~Perron and V.~Furnon.
\newblock OR-Tools (CP-SAT solver).
\newblock Google, \url{https://developers.google.com/optimization}.

\bibitem{PM26}
M.~Perez and J.~Flores.
\newblock The strict simplicial complexity of the vertex-minimal
  M\"obius band triangulation is~$2$.
\newblock Companion note, 2026.

\end{thebibliography}

\end{document}





